\documentclass[12pt,final,twoside]{report}
%%%%%%%%%%%%%%%%%%%%%%%%%%%%%%%%%%%%%%%%%%%%%%%%%%%%%%%%%%%%%
% Some credits:
% The template initially was created by Prof. Dr. Holger Karl/Uni Paderborn '2006
% and was expanded and updated by Dipl.-Inform. Stefan Heinrich/Uni Paderborn/Uni Hamburg since 2008.
% Suggestions for changes are always welcome.
%%%%%%%%%%%%%%%%%%%%%%%%%%%%%%%%%%%%%%%%%%%%%%%%%%%%%%%%%%%%%
% Meta information:
\newcommand{\trtitle}{Low-Dimensional Control under Embodiment Constraints: From LQR to Reinforcement Learning}
\newcommand{\trtype}{Bachelor's Thesis}
\newcommand{\trauthor}{Alexander Michael Kranz}
\newcommand{\trauthortitle}{} %{Dipl.-Inform.\ }
\newcommand{\tremail}{alexander.kranz@tuhh.de}
\newcommand{\trmatrikelnummer}{590156}
\newcommand{\trbirth}{09.02.2004}
\newcommand{\trstudies}{Data Science}
\newcommand{\trgutachterA}{\href{mailto:nihat.ay@tuhh.de}{Prof. Dr. N. Ay}}
\newcommand{\trgutachterB}{\href{mailto:manfred.eppe@tuhh.de}{Dr. Manfred Eppe}}
\newcommand{\trbetreuung}{\href{mailto:adwait.datar@tuhh.de}{Dr. Adwait Datar}}
\newcommand{\trfach}{Data Science Foundations}
\newcommand{\trdate}{31.10.2026}
\newcommand{\trkeywords}{LQR, Cheap Design, Reinforcement Learning, Low-Dimensional Control, Embodied Agents}

%%%%%%%%%%%%%%%%%%%%%%%%%%%%%%%%%%%%%%%%%%%%%%%%%%%%%%%%%%%%%
% Languages:

% Falls die Ausarbeitung in Deutsch erfolgt:
% \usepackage[german]{babel}
% \usepackage[T1]{fontenc}
% \usepackage[latin1]{inputenc}
% \usepackage[latin9]{inputenc}                     
% \selectlanguage{german}

% If the thesis is written in English:
\usepackage[english]{babel}                         
\selectlanguage{english}

%%%%%%%%%%%%%%%%%%%%%%%%%%%%%%%%%%%%%%%%%%%%%%%%%%%%%%%%%%%%%
% Bind packages:
\usepackage{acronym}                    % Acronyms
\usepackage{algorithmic}                % Algorithms and Pseudocode
\usepackage{algorithm}                  % Algorithms and Pseudocode
\usepackage{amsfonts}                   % AMS Math Packet (Fonts)
\usepackage{amsmath}                    % AMS Math Packet
\usepackage{amssymb}                    % Additional mathematical symbols
\usepackage{amsthm}
\usepackage{booktabs}                   % Nicer tables
%\usepackage[font=small,labelfont=bf]{caption} % Numbered captions for figures
\usepackage{color}                      % Enables defining of colours via \definecolor
\definecolor{uhhRed}{RGB}{226,0,26}     % Official Uni Hamburg Red
\definecolor{uhhGrey}{RGB}{136,136,136} % Official Uni Hamburg Grey
\definecolor{uhhLightGrey}{RGB}{220, 220, 220}
\usepackage{fancybox}                   % Put equations in a frame
\usepackage{fancyhdr}                   % Packet for nicer headers
%\usepackage{fancyheadings}             % Nicer numbering of headlines
\usepackage[body={5.8in,9in}]{geometry} % Type area (size, margins...)

\geometry{a4paper,outer=3.35cm}        % !!!Release version (Normal margins)
%\geometry{a4paper,outer=2.5cm}         % !!!Print version (Additional margin on the left for the binding)
%\geometry{a4paper}                     % !!!Proofread version (Additional margin on the right for corrections)
%\geometry{a4paper,outer=3.15cm}         % !!!Draft version (Same margins on left and right)
%\geometry{paperheight=10.0in,paperwidth=6.4in,top=0.51in,left=0.3in}  % !!!Developer version (Minimal margins)

\usepackage{graphicx}                   % Inclusion of graphics
%\usepackage{latexsym}                  % Special symbols
\usepackage{longtable}                  % Allow tables over several parges
\usepackage{listings}                   % Nicer source code listings
\usepackage{multicol}                   % Content of a table over several columns
\usepackage{multirow}                   % Content of a table over several rows
\usepackage{rotating}                   % Alows to rotate text and objects
% \usepackage[hang]{subfigure}            % Allows to use multiple (partial) figures in a fig
%HINT: subfigure may be outdated already - maybe try subfig instead 
%\usepackage[font=footnotesize,labelfont=rm]{subfig}    % Pictures in a floating environment
\usepackage{tabularx}                   % Tables with fixed width but variable rows
\usepackage{caption}
\usepackage{subcaption}
\usepackage{url,xspace,boxedminipage}   % Accurate display of URLs
\usepackage{parskip} % blank line between paragraphs, no indent

%%%%%%%%%%%%%%%%%%%%%%%%%%%%%%%%%%%%%%%%%%%%%%%%%%%%%%%%%%%%%
% PDF Information und Definitions:
\author{\trauthor}

\ifx\pdftexversion\undefined
\usepackage{hyperref}
\else
\usepackage[colorlinks=false,           % link is colores (true) or has colored frame (false)
            linkcolor=uhhRed,           % case colorlinks=true: define color.
            urlcolor=uhhRed,
            citecolor=uhhRed,
            bookmarks,                  % Place bookmarks erstellen
            bookmarksopen=true,         % Bookmarks will be shown at start (true/false)
            pdfpagemode=UseOutlines,    
            bookmarksopenlevel=1,       % Define the depth of shown links
            bookmarksnumbered,          % Numbers of chapers in Bookmarks
            pdftitle={\trtitle},
            pdfsubject={\trtype},
            pdfkeywords={\trkeywords},
            pdfauthor={\trauthor},
            plainpages=false
            ]{hyperref}
\fi

\ifx\pdftexversion\undefined
\else
\pdfoutput=1                            % Disable PDF-Output
\pdfimageresolution=1200
\pdfcompresslevel=2                     % 0 = no compression, 9 = strongest compression
\fi

%%%%%%%%%%%%%%%%%%%%%%%%%%%%%%%%%%%%%%%%%%%%%%

% Self-included packages, not from base
\usepackage{amsmath, amssymb, mathtools}
\usepackage{pgf}
\usepackage{tikz}
\usetikzlibrary{arrows.meta,positioning}
\usepackage{import}
\usepackage{xcolor}
\usepackage[capitalise,nameinlink]{cleveref}
\usepackage{placeins}                   % \FloatBarrier: keep floats within their section

%%%%%%%%%%%%%%%%%%%%%%%%%%%%%%%%%%%%%%%%%%%%%%%%%%%%%%%%%%%%%
% Configurationen:

\hyphenation{whe-ther}                  % Manually use: "\-" in a word: Staats\-ver\-trag

%\lstloadlanguages{C}                   % Set the default language for listings
\DeclareGraphicsExtensions{.pdf,.svg,.jpg,.png,.eps} % first try pdf, then eps, png and jpg
\graphicspath{{./src/}}                 % Path to a folder where all pictures are located
\pagestyle{fancy}                       % Use nicer header and footer

% Redefine the environments for floating objects:
\setcounter{topnumber}{3}
\setcounter{bottomnumber}{2}
\setcounter{totalnumber}{4}
\renewcommand{\topfraction}{0.9}        %Standard: 0.7
\renewcommand{\bottomfraction}{0.5}     %Standard: 0.3
\renewcommand{\textfraction}{0.1}       %Standard: 0.2
\renewcommand{\floatpagefraction}{0.8}  %Standard: 0.5

% Tables with a nicer padding:
\renewcommand{\arraystretch}{1.2}

\setcounter{secnumdepth}{3}

% Chapter and Sections will not be written in capitals
\renewcommand{\chaptermark}[1]{\markboth{\chaptername \ \thechapter.\ #1}{}}
\renewcommand{\sectionmark}[1]{\markright{\thesection.\ #1}}

%%%%%%%%%%%%%%%%%%%%%%%%%%%%
% Additional 'theorem' and 'definition' blocks:
\theoremstyle{plain}
\newtheorem{theorem}{Theorem}[chapter]
%\newtheorem{theorem}{Satz}[chapter]    % Wenn in Deutsch geschrieben wird.
\newtheorem{axiom}{Axiom}[chapter]     
%\newtheorem{axiom}{Fakt}[chapter]      % Wenn in Deutsch geschrieben wird.
%Usage:%\begin{axiom}[optional description]%Main part%\end{fakt}

\theoremstyle{definition}
\newtheorem{definition}{Definition}[chapter]

%Additional types of axioms:
\newtheorem{lemma}{Lemma}[chapter]
\newtheorem{observation}[axiom]{Observation}

%Additional types of definitions:
\theoremstyle{remark}
%\newtheorem{remark}[definition]{Bemerkung} % Wenn in Deutsch geschrieben wird.
\newtheorem{remark}[definition]{Remark} 

%%%%%%%%%%%%%%%%%%%%%%%%%%%%
% Provides TODOs within the margin:
\newcommand{\TODO}[1]{\marginpar{\emph{\small{{\bf TODO: } #1}}}}

%%%%%%%%%%%%%%%%%%%%%%%%%%%%%%%
% Own commands

% To import a matplotlib figure
\newcommand{\generatedfigure}[2]{\import{generated/#1/figures/}{#2.pgf}}
% Half-width figures
\newcommand{\generatedhalffigure}[2]{\import{generated/#1/figures/half/}{#2.pgf}}

% Needs citations command
\newcommand{\needcite}[1]{\textcolor{red}{\textsuperscript{[CITE: #1]}}}

% Macros for notation
\newcommand{\R}{\mathbb{R}}
\newcommand{\C}{\mathbb{C}}
\newcommand{\E}{\mathbb{E}}
\newcommand{\T}{^{\top}}
\newcommand{\Htwo}{\mathcal{H}_{2}}
\newcommand{\Hinf}{\mathcal{H}_{\infty}}
\DeclareMathOperator{\tr}{tr}
\DeclareMathOperator{\rank}{rank}
\DeclareMathOperator{\spr}{\rho}

%%%%%%%%%%%%%%%%%%%%%%%%%%%%%%%%%%%%%%%%%%%%%%%%%%%%%%%%%%%%%
% Document:

\begin{document}

\pagenumbering{Roman}                   % Roman pagenumbering for lists and meta pages
\renewcommand{\headheight}{14.5pt}      % Size of headings

\thispagestyle{empty}
\fancyhead[LO,RE]{}                     % Define the header style for the meta pages

%%%%%%%%%%%%%%%%%%%%%%%%%%%%
% Cover sheet

\begin{titlepage}
%---Possibility 1:
    \begin{flushleft}
        % trim={left bottom right top}
        \includegraphics[width=65mm,trim={5cm 14cm 0 13cm},clip]{TUHH_logo_rgb.pdf} \hfill \includegraphics[width=45mm,trim={0 0.5cm 0 0cm},clip]{dsf-logo-2024.png}\\[10pt]
        Technische Universität Hamburg \hfill Data Science Foundations
    \end{flushleft}

    \rule{\textwidth}{0.4pt}
        \newline
        \vspace{2.0cm}
        \begin{center}
          \LARGE \textbf{\trtitle}
        \end{center}
    \vspace{2.0cm}
    \begin{center}
      \textbf{\trtype}\\
      am Institut \trfach\\
      Technische Universität Hamburg \\[0.5cm]
      vorgelegt von \\
      \textbf{\trauthortitle\href{mailto:\tremail}{\trauthor}}\\
      am\\
      \trdate
    \end{center}
    \vspace{1cm}
    \begin{center}
    \begin{tabular}{ll}
    Erstprüfer: & \trgutachterA \\
    Zweitprüfer & \trgutachterB \\
    Betreuung:  & \trbetreuung \\    	% Adviser are not allowed to demand getting credited, but are happy getting credited by the students initiative
    \end{tabular}
    \end{center}
    \vfill
    \begin{tabular}{l}
    \trauthor \\
    geboren am: \trbirth \\
    Matrikelnummer:  \trmatrikelnummer \\
    Studiengang: \trstudies
    \end{tabular}
    \newline
    \rule{\textwidth}{0.4pt}
    \newpage 
\end{titlepage}

    %backsite of cover sheet is empty!
\thispagestyle{empty}
\hspace{1cm}
\newpage

%%%%%%%%%%%%%%%%%%%%%%%%%%%%
% Abstract:
\section*{Abstract}
\label{sec:abstract}
\fancyhead[LE,RO]{\it Abstract}
%\addcontentsline{toc}{chapter}{\numberline{}Abstract}
%Your English abstract here (mandatory if written in English and recommended otherwise).
This thesis investigates the effects of low-dimensional structures in the control architectures for
embodied agents, starting from simple LQR-based solvers and expanding to RL-approaches implemented using proximal
policy optimization. As the basis serves the double integrator, which is first artificially overparameterized by
having the controller observe and act in a higher dimension and use projection matrices to change between it and the
dimension of the underlying system. Afterwards, we also introduce a system with camera-based observations and a more complex actuation structure.
Throughout our experiments, we notice that low-dimensional controllers are difficult to implement successfully, and
require careful consideration during design. However, if done right, a combination of fixed, briefly pretrained encoders
together with a smaller, longer-trained actor can perform similarly to a fully connected one. With regard to noise,
no significant improvement in robustness for low-dimensional controllers could be observed.
However, a fully trained controller seems to be more robust to control noise.

% If the abstract ist not longer than half a page, then the German Zusammenfassung can be places on the same page
%\cleardoublepage
\vfill
\section*{Zusammenfassung}
\label{sec:zusammenfassung}
%\addcontentsline{toc}{chapter}{\numberline{}Zusammenfassung} %Add the Zusammenfassung to the TOC
%Hier die deutsche Zusammenfassung einf\"ugen (notwendig).
%\fancyhead[LE,RO]{\it Abstract}
Diese Arbeit untersucht die Effekte von niederdimensionalen Strukturen in der Regelungsarchitektur von
Agenten mit einem Körper. Hierfür beginnen wir mit LQR-basierten Reglern und erweitern unsere Ansätze
danach zu Reinforcement Learning, für welches wir Proximal Policy Optimization verwenden. Die Basis
bildet der Doppelintegrator, welchen wir künstlich überparametrisieren, indem wir den Regler in einer
höheren Dimension beobachten und agieren lassen, und mit Hilfe von Projektionsmatrizen zwischen dieser
und der echten Dimension des Systems transferieren. Danach stellen wir ein kamera-basiertes System vor,
welches eine komplexere Motorenstruktur besitzt. In unseren Experimenten erkennen wir, dass es schwierig
ist niederdimensionale Regler effektiv zu implementieren und in jedem Falle genaue Planung bei dem Aufbau
benötigt. Bei korrekter Umsetzung können Kombinationen aus statischen, vortrainierten Codierern und kleineren,
länger trainierten neuronalen Netzen jedoch Leistungen erzielen, die vergleichbar mit einem ganzheitlich trainierten
Netz sind. In Bezug auf Rauschen konnten keine signifikanten Verbesserungen in der Resistenz für niederdimensionale Regler
beobachtet werden. Allerdings scheint ein ganzheitlich trainierter Regler resistenter
gegenüber Regelungsstörungen zu sein.

\cleardoublepage

%%%%%%%%%%%%%%%%%%%%%%%%%%%%
% Acknowledgements:
\section*{Acknowledgements}
\label{sec:acknowledgements}
\fancyhead[LE,RO]{\it Acknowledgements}
I want to thank Dr.~Adwait Datar for his guidance and expertise, as well as my friends and family for their
continuous support during my thesis.

\cleardoublepage

%%%%%%%%%%%%%%%%%%%%%%%%%%%%
% Lists:
%\setcounter{tocdepth}{1}               % depth of the table of contents (for BSc and MSc Thesis 1 is recommented)
\fancyhead[LE,RO]{\it Contents}
\tableofcontents
\cleardoublepage
% List of Figures and List of tables are optional. -> Not needed in most theses.
\fancyhead[LE,RO]{\it List of Figures}
\listoffigures
\cleardoublepage
\fancyhead[LE,RO]{\it List of Tables}
\listoftables
\cleardoublepage
%\lstlistoflistings
%\cleardoublepage

\fancyhead[LE]{\it \leftmark}           % Define the header style for the text pages
\fancyhead[RO]{\it \rightmark}          % Define the header style for the text pages
\fancyhead[LO,RE]{}                     % Define the header style for the text pages

%%%%%%%%%%%%%%%%%%%%%%%%%%%%
% The content will be included here:
\pagenumbering{arabic}

\chapter{Introduction}\label{introduction}


% The introduction describes the problem, why it is important, the main
% ideas of the thesis, what are the main contributions of the
% thesis, etc. The introduction is the most important part of your thesis because it answers the question of why you are researching your topic at all. The introduction does not only motivate your thesis, it also justifies it. 
% Compared to the rest of your thesis, it is a very short text, but this text should be extremely well-written and clear. It deserves a lot of polishing. If the introduction is boring or not well written, a reader (and your supervisor who will finally grade your thesis) will most likely not have fun reading your work. The introduction (together with abstract and title) is the ``bait'' to get your reader, your supervisor and your reviewer hooked.

% The introduction should be a short version of the whole thesis, with a focus on motivation. A good introduction is structured as follows:

% \begin{enumerate}
% \item It should start with a brief motivational problem outline to set the stage, e.g., ``It is important to build intelligent robots. A popular method to build intelligent robots is X. However, X has problem Y.''

% \item It should then contain a specific formulation of your research question (RQ). For a good description of how to formulate a research question, consider \url{https://www.scribbr.com/research-process/research-questions/}. E.g.: ``Problem Y has not yet been addressed sufficiently by other researchers, so how can we extend X to address Y?''

% \item In addition to your research question, you can (but don't have to) state a specific hypothesis: 
% ``I hypothesize that we can address Y by extending X with Z.'' It is important that you also say how you plan to validate your hypothesis. For example, if you do an empirical investigation, write something like ``I will validate my hypothesis by measuring ... with experiment ...''. If you write a more formal thesis, you might want to prove a certain theorem, so you could write ``I will validate my hypothesis by proving that ...''

% \item The RQ and hypothesis should be followed by a very brief (one-two sentences) summary of how you want to address the question, i.e., the methods and evaluation metrics you want to use, highlighting your research contributions with respect to the state of the art. 
% Herein, you should also refer to the individual sections in your paper to give the reader a brief outline. 
% ``The review of the related work in Section 2 of this paper shows that method Q has not yet been investigated as a solution to problem X. As a novel contribution, I will address this hypothesis / RQ by implementing Z using method Q, as described in Section 3. I will evaluate the results by measuring V, as described in Section 4.'' (Note that you might not have a Section 4 if you do not conduct empirical research  but focus on theory instead.)

% \item You should conclude with a sentence like this: ``I conclude in Section 5, by highlighting that extending X with Z indeed achieves Y, but only if W holds.'' 

% \end{enumerate}
% It does not matter whether you use ``I'' or ``We''. It does matter, however, that you use the same pronoun consistently.

% My part
Undoubtedly, the field of robotics has seen rapid advancements in the capabilities of neural network based implementations over the recent years.
An increase in global compute power~\cite{TODO_compute}, fueled by the so-called ``AI race'' of LLMs, has allowed researchers around the world to create even more extensive model architectures for more complex use cases~\cite{TODO_llm_example}.
Transferring these structures now to physical bodies, however, shows the limits of this approach, and prompts the idea of cheap design.
As mentioned by N.~Ay~\cite{TODO_ay_geometric}, the core idea is to maintain an ``embodied universal approximator'', while reducing the overall complexity of the model, which can be defined based on the specific case.
The most straightforward approach is inducing the underlying physical dynamics of the embodiment into the architecture, which is often of lower-dimensionality than the actual sensors and actuators.
Recent papers focusing on algorithms for the LQR~\cite{TODO_lqr}, apply this approach by offering implementations for enforcing low-rank controllers, leading to lower complexity in relation to parameterization.
Still, this approach has yet to be extended~\cite{TODO_extension} to more complex embodiments, and is missing a crucial part required for implementation: Robustness.
This leaves us with the question: \emph{Is it possible to train a controller of low-dimensionality for an embodied agent, which is performant yet robust?}

To answer it, we will first explain the core concepts and collect existing approaches for the LQR setting in section~\ref{preliminaries} and \ref{background}.
In section~\ref{lqr}, we will compare these implementations using the double integrator problem for LQR with respect to general performance, as well as robustness.
Building on that, we will artificially inflate our inputs and outputs via matrix-multiplication in section~\ref{inflation} and introduce additionally stable-baseline models into the comparison.
For section~\ref{perception}, we will completely move away from LQR, and extend the setting to a perception-based implementation.
Here we will compare different models with varying parameter complexity in regards to robustness, applying the concepts we learned from our earlier experiments.
Finally, we will conclude by analyzing whether we indeed achieved our goal of constructing a robust controller, and stating how our approach can be applied to more complex settings.
\cleardoublepage

\chapter{Background}
\label{chap:background}

Before turning to the experiments, we need to lay the groundwork and
reasoning for them. We start by defining some important concepts from
control theory, as well as introducing the system that serves as our
core example throughout (\cref{sec:preliminaries}). Then, we discuss
the current state of research around embodiment, low-dimensional
control, and in that context robustness as well (\cref{sec:motivation}).
We further reason about the relevance of this thesis, as well as comparable
papers. After that, we move on to introducing our specifically chosen
setup for our experiments, define constants used and explain
our test protocol (\cref{sec:setup}). Finally, we walk the reader
through our upcoming chapters, building on the previously introduced
information, and giving reason for certain design choices (\cref{sec:approach}).

\section{Preliminaries}
\label{sec:preliminaries}

Here, we revisit the concepts that are necessary for understanding the
following chapters. This section introduces basics of control theory, such as
linear state-space models (\cref{subsec:state-space}), the double
integrator as our example (\cref{subsec:double-integrator}), and the
linear quadratic regulator (\cref{subsec:lqr}).

\subsection{Linear State-Space Models}
\label{subsec:state-space}
 
In control theory, the evolution of an agent together with its environment is
described by a \emph{state-space model}: a rule that assigns to every point of
a state space the instantaneous rate at which the system moves through it.
Throughout this thesis we restrict attention to the class of
\emph{linear time-invariant} systems, for which this rule is linear in
both the state and the input and does not change over time~\cite[pp.~3--4]{Werner2021Control}.

Such a continuous-time system is given by
\begin{subequations}
\label{eq:ss-cont}
\begin{align}
    \dot{x}(t) &= A_c\, x(t) + B_c\, u(t), \qquad x(0) = x_0, \label{eq:ss-cont-state}\\
    y(t)       &= C\, x(t) + D\, u(t), \label{eq:ss-cont-output}
\end{align}
\end{subequations}
where $x(t) \in \R^{n}$ is the \emph{state vector}, $u(t) \in \R^{m}$ the
\emph{input vector} chosen by the agent, and $y(t) \in \R^{q}$ the
\emph{output vector}, i.e.\ the output observed and thus available to the
feedback controller. Accordingly, $A_c \in \R^{n \times n}$ is called the
\emph{system matrix}, $B_c \in \R^{n \times m}$ the \emph{input matrix},
$C \in \R^{q \times n}$ the \emph{output matrix}, and $D \in \R^{q \times m}$
the \emph{feedthrough matrix}.
 
The two equations in \eqref{eq:ss-cont} play fundamentally different roles, and
it is helpful to explain their distinction.
\Cref{eq:ss-cont-state} is a differential equation: it does not compute the
state, which is assumed to be known at every instant, but prescribes its
\emph{rate of change}. The pair $(A_c, B_c)$ therefore governs how the system
\emph{behaves}, with $A_c$ describing its autonomous dynamics and
$B_c$ describing how the input of the agent influences them. \Cref{eq:ss-cont-output}, in contrast, 
contains no derivative at all. It is an instantaneous readout, and the pair $(C, D)$ therefore
describes how the system is \emph{observed} rather than how it evolves. We set
$D = 0$ throughout, which is the standard assumption for systems in which the
actuator has no direct instantaneous influence on the sensor~\cite[p.~4]{Werner2021Control}.
 
Given an input signal $u(\cdot)$, the trajectory of \eqref{eq:ss-cont-state} is
obtained by integration and is given in closed form by the variation-of-constants
formula~\cite[p.~13, Eq.~(2.20)]{Werner2021Control}
\begin{equation}
    \label{eq:variation-of-constants}
    x(t) = e^{A_c t} x_0 + \int_{0}^{t} e^{A_c (t - s)} B_c\, u(s)\, \mathrm{d}s ,
\end{equation}
where $e^{A_c t}$ denotes the matrix exponential.
 
\paragraph{Full-state measurement.}
Until \cref{subsubsec:inflated-model}, we assume that the state is fully
measurable and therefore set $C = I_n$, which reduces
\eqref{eq:ss-cont-output} to $y(t) = x(t)$. Later,
we replace $C$ by a lift matrix $U_1$ and
inject measurement noise (\cref{subsubsec:noise-injection}),
such that the previous simplification no longer holds.
 
% -----------------------------------------------------------------------------
\subsubsection{Discretization}
\label{subsubsec:discretization}
 
Digital controllers, as in our case, do not act continuously: they sample the output at a fixed
rate and hold the computed control input constant until the next sampling point. The
learning algorithms studied in this thesis operate on such sampled trajectories
as well, so our model has to reflect this. The standard approach is the
\emph{zero-order hold} (ZOH)~\cite[p.~118]{Werner2021Control}. Fixing a sampling period $h > 0$ and writing
$x_k \coloneqq x(kh)$, the input is assumed to be piecewise constant~\cite[Eq.~(7.1)]{Werner2021Control},
\begin{equation}
    \label{eq:zoh-input}
    u(t) = u_k \qquad \text{for } t \in [kh,\, (k+1)h), \quad k \in \mathbb{N}_0 .
\end{equation}
Substituting \eqref{eq:zoh-input} into \eqref{eq:variation-of-constants} over a
single sampling interval and applying the change of variables
$\tau = (k+1)h - s$ yields the discrete-time model
\begin{equation}
    \label{eq:ss-disc}
    x_{k+1} = A_d\, x_k + B_d\, u_k ,
\end{equation}
with~\cite[p.~136, Eqs.~(7.26)--(7.28)]{Werner2021Control}
\begin{equation}
    \label{eq:zoh-matrices}
    A_d = e^{A_c h},
    \qquad
    B_d = \left( \int_{0}^{h} e^{A_c \tau} \, \mathrm{d}\tau \right) B_c .
\end{equation}

We use the implementation from SciPy~\cite{SciPy2026ContDiscrete}
to carry out this transformation in the code. Unless stated otherwise, all experiments in this
thesis are carried out in discrete time with sampling period $h = 0.01\,\mathrm{s}$, and
we drop the subscript $d$ in the following, writing simply $A$ and $B$.
 
\subsection{The Double Integrator}
\label{subsec:double-integrator}
 
To make the comparisons across all chapters
meaningful, all experiments are built on a single system: the
\emph{double integrator}, a frictionless point mass of unit mass moving along a
line under an applied force $u$~\cite[pp.~87--88]{Rao2001Naive}. Its dynamics are Newton's second law,
\begin{equation}
    \label{eq:newton}
    \ddot{p}(t) = u(t),
\end{equation}
where $p(t)$ denotes the position of the mass. We simplify by using a mass of one, so that
force and acceleration coincide numerically. Choosing position and velocity as
the state, $x = [\,p,\ \dot{p}\,]\T \in \R^{2}$, the second-order equation
\eqref{eq:newton} decomposes into the first-order system
\begin{equation}
    \label{eq:di-first-order}
    \frac{\mathrm{d}}{\mathrm{d}t} \begin{bmatrix} p(t) \\ \dot{p}(t) \end{bmatrix}
    = \begin{bmatrix} \dot{p}(t) \\ u(t) \end{bmatrix},
\end{equation}
in which the first row is the definition of velocity and the second is
\eqref{eq:newton} itself. In the notation of \eqref{eq:ss-cont} this gives
\begin{equation}
    \label{eq:di-cont-matrices}
    A_c = \begin{bmatrix} 0 & 1 \\ 0 & 0 \end{bmatrix},
    \qquad
    B_c = \begin{bmatrix} 0 \\ 1 \end{bmatrix},
    \qquad
    C = I_2, \qquad D = 0 .
\end{equation}
The entry $(A_c)_{12} = 1$ should be read as stated in \cref{subsec:state-space}:
the rate of change of position is given by the current velocity.
 
Because $A_c^{2} = 0$, the matrix exponential in
\eqref{eq:zoh-matrices} truncates after two terms, $e^{A_c h} = I + A_c h$, and
the discretization is available in closed form:
\begin{equation}
    \label{eq:di-disc-matrices}
    A_d = \begin{bmatrix} 1 & h \\ 0 & 1 \end{bmatrix},
    \qquad
    B_d = \begin{bmatrix} h^{2}/2 \\ h \end{bmatrix}.
\end{equation}
This system is the minimal nontrivial choice for our purposes. It is
second-order, so the state is not directly actuated and the controller must
account for momentum. It is \emph{controllable} for every $h > 0$, i.e.\ an input
sequence can navigate it from any initial state to any final
state~\cite[p.~127, Def.~7.2]{Werner2021Control}.%
\footnote{Controllability, with its weaker form stabilizability, and the dual notions
observability and detectability are standard structural properties of linear systems.
If the reader wishes to read up on their formal definitions, we kindly refer them
to~\cite[Chs.~3--4 and pp.~127--129]{Werner2021Control}.}
Additionally, it is not open-loop stable: an unforced discrete-time system $x_{k+1} = A x_k$ is
\emph{stable} if and only if all eigenvalues of $A$ lie strictly inside the unit
disk~\cite[p.~127, Thm.~7.1]{Werner2021Control}, i.e.\ if its \emph{spectral radius}
$\spr(A) \coloneqq \max_i \lvert \lambda_i(A) \rvert$ satisfies $\spr(A) < 1$. Here, $A_c$
has a double eigenvalue at the origin and $A_d$ correspondingly a double eigenvalue at $1$,
so $\spr(A_d) = 1$. The last property matters for the
learning experiments: an uncontrolled or badly initialized policy diverges, so
the assumption of a known stabilizing initial gain (\cref{subsec:protocol}) is a
restriction.
 
\subsection{The Linear Quadratic Regulator}
\label{subsec:lqr}

Now that we have defined our state-space model and the system we
study, it is time to introduce our controller, which is based on the
\emph{linear quadratic regulator} problem.

Here, the setting is that our system starts at an arbitrary initial state in its
state space, and the goal of our controller is to return to the zero state as
cheaply as possible. Both distance to the origin, as well as applied control
input yield certain costs, which can be combined into a single cost function, where $Q$
describes the former, while $R$
associates costs to control input. Over a finite horizon $T$, this gives the
cost~\cite[pp.~5--6, Eqs.~(1.3) and (1.5)]{Werner2021Optimal}
\begin{equation}
    \label{eq:lqr-cost-finite}
    V_T(x_0, u) = \int_{0}^{T} \left( x(t)\T Q\, x(t) + u(t)\T R\, u(t) \right) \mathrm{d}t ,
    \qquad Q = Q\T \succeq 0, \quad R = R\T \succ 0 .
\end{equation}
We notice two further restrictions here: $Q$ is required to be positive
semi-definite, meaning that deviating from the origin never reduces the cost,
and $R$ must be positive definite, such that every control input results in a cost.

As the original formula relates to a continuous problem with a fixed end time $T$,
it does not fit well with our goal of optimizing over an infinite discrete horizon.
The following formula offers exactly that~\cite[p.~137, Eq.~(7.30)]{Werner2021Control}:
\begin{equation}
    \label{eq:lqr-cost-infinite}
    J(x_0, u) = \sum_{k=0}^{\infty} \left( x_k\T Q\, x_k + u_k\T R\, u_k \right) .
\end{equation}
This sum does not necessarily converge. If the closed loop
does not decay, the costs do not vanish and $J$ tends to $+\infty$. 
Using this as our definition, $J$ becomes well defined for every gain.

\subsubsection{The Optimal Gain}
\label{subsubsec:lqr-optimal}

What makes the LQR problem special is that its solution is a static linear feedback
law~\cite[Sec.~4]{Recht2019Tour}: the optimal input at time $k$ depends only on the current state of the model.
Introducing a gain matrix $K \in \R^{m \times n}$, we write it as%
\footnote{Note that the minus in front of $K$ has been chosen as the convention for this thesis.
This turns our closed loop into the format $A - BK$, while some referenced literature
uses $A + BK$. Formulas cited have been adapted to our convention by replacing $K$ with $-K$.
}
\begin{equation}
    \label{eq:feedback-law}
    u_k = -K x_k .
\end{equation}
Our controller is thus reduced to choosing a matrix, and the
optimal controller is given by the optimal gain matrix $K^{*}$.

As we are dealing with a discrete-time system, that gain is obtained from the
\emph{discrete algebraic Riccati equation} (DARE)~\cite[Sec.~4, Eq.~(4.2)]{Recht2019Tour}
\begin{equation}
    \label{eq:dare}
    P = A\T P A - A\T P B \left( R + B\T P B \right)^{-1} B\T P A + Q ,
\end{equation}
whose stabilizing solution $P = P\T \succeq 0$ yields~\cite[Sec.~4]{Recht2019Tour}
\begin{equation}
    \label{eq:optimal-gain}
    K^{*} = \left( R + B\T P B \right)^{-1} B\T P A ,
    \qquad
    u_k^{*} = -K^{*} x_k .
\end{equation}
The matrix $P$ is additionally our
\emph{cost-to-go} matrix~\cite[Sec.~1, Eq.~(1)]{Fazel2018Global}, meaning that we can use it to calculate the optimal cost 
incurred from any state $x$ onwards with
\begin{equation}
    \label{eq:cost-to-go}
    J^{*}(x) = x\T P x .
\end{equation}

Whether \eqref{eq:dare} has such a solution at all depends on two structural
properties of the system, stabilizability and
detectability~\cite[p.~442, Thm.~8]{Kucera1972Discrete}:

\begin{theorem}[Existence and uniqueness]
\label{thm:dare-existence}
Let $(A, B)$ be stabilizable and let $(A, Q^{1/2})$ be detectable. Then the DARE
\eqref{eq:dare} has a unique symmetric positive semi-definite solution $P$
and the resulting closed loop satisfies $\spr(A - B K^{*}) < 1$.
\end{theorem}

For the double integrator, both conditions hold: $(A, B)$ is controllable and thus
stabilizable, and with $Q = I_n$ (\cref{subsec:protocol}) the pair $(A, Q^{1/2}) = (A, I_n)$
is also observable. Thus $K^{*}$ is well defined in all our experiments.

\subsubsection{The Cost as a Function of the Gain}
\label{subsubsec:lqr-cost-of-gain}

\Cref{eq:optimal-gain} solves the problem, but it solves it directly and
tells us nothing about the quality of non-optimal gains. As the learning 
algorithms of the following chapters will produce such $K$ directly, 
we need to be able to calculate the cost as a function of $K$.

Fix a gain $K$ and write the resulting closed loop as
$x_{k+1} = A_K x_k$ with $A_K \coloneqq A - BK$. Substituting
\eqref{eq:feedback-law} into \eqref{eq:lqr-cost-infinite} shows that the cost is
again a quadratic form, $J(x_0, K) = x_0\T P_K x_0$, where $P_K$ now solves
the far simpler discrete \emph{Lyapunov equation}~\cite[Sec.~2.1]{Fazel2018Global}
\begin{equation}
    \label{eq:lyapunov}
    P_K = A_K\T P_K A_K + Q + K\T R K ,
\end{equation}
which is linear in the unknown $P_K$ and has a unique solution whenever
$\spr(A_K) < 1$. This restricts us to the set of stabilizing gains
\begin{equation}
    \label{eq:stabilizing-set}
    \mathcal{K} \coloneqq \left\{ K \in \R^{m \times n} : \spr(A - BK) < 1 \right\} ,
\end{equation}
outside of which we set the cost to $+\infty$, consistent with the remark after
\eqref{eq:lqr-cost-infinite}.

One dependency remains, namely on the initial state. Evaluating the gain
based on a single $x_0$ would reward gains that happen to fit that one
point, so we instead draw the initial state from a distribution,
$x_0 \sim \mathcal{D}$, and average~\cite[Sec.~3]{Fazel2018Global}. Writing
$\Sigma_0 \coloneqq \E_{x_0 \sim \mathcal{D}}[x_0 x_0\T]$ for its second moment,
the objective~\cite[Sec.~2.1]{Fazel2018Global} becomes
\begin{equation}
    \label{eq:cost-of-gain}
    f(K) \coloneqq \E_{x_0 \sim \mathcal{D}} \left[ x_0\T P_K x_0 \right]
    = \tr \left( P_K \Sigma_0 \right) ,
\end{equation}
using the cyclic property of the trace. \Cref{eq:cost-of-gain} is the objective
that the algorithms in \cref{chap:lqr} minimize and they differ
only in how much of the model is known to them. Later, in \cref{chap:inflation,chap:perception},
the RL approaches rely on a rollout-based estimate of it to optimize over.

% ------------------------------------------

\section{Motivation}
\label{sec:motivation}

An agent acting in its environment receives observational input and responds to it with a commanded
control input.
In other words, the agent is a function in relation to its environment, and to find an optimal
behavior, we need to approximate this optimal function as closely as possible. While neural networks
are known to be universal approximators~\cite{Cybenko1989Approximation,Hornik1989Multilayer}, and
hence able to approximate any function, in the case of an embodied agent, different constraints need to
be taken into account. Ay~\cite[p.~392]{Ay2015Geometric} names additionally to external embodiment constraints
possible intrinsic constraints, and argues that such can even be ``beneficial'' for the control
structure. By the core notion of cheap design, embodiment can help produce complex
behaviors while only needing smaller controllers~\cite{Pfeifer2006HowBody}, which
inherently requires constraints. So, for such embodied agents, the goal is to find an ``embodied universal
[approximator]''~\cite[p.~3]{Montufar2015CheapControl}, which restricts the neural network
architecture to an intrinsic, unknown dimension, encoding physical properties belonging to the agent
and the environment, while disallowing impossible movements~\cite[Sec.~3]{Montufar2015CheapControl}.

Biology motivates this approach even further: here it is clear to see that by commanding a body,
some features of the control architecture can be simplified and delegated to the
embodiment~\cite{Mueller2017Morphological}. For example, McGeer's passive walker~\cite{McGeer1990Passive} clearly
shows that some movements can be produced by the body without any control input, just by choosing
the right embodiment. While obviously here being simplified to not include any control at all, it
can quite simply be extended with such~\cite[Sec.~10]{McGeer1988Stability}.
M{\"u}ller and Hoffmann~\cite[p.~4]{Mueller2017Morphological} mention such a combination, a robot
climbing up walls, whose design was inspired by geckos, indeed using energy for control, however
very little compared to an approach not exploiting this specific embodiment design. Further, it can
be seen that the complex
behaviors are not only less energy-intensive to produce~\cite[p.~4]{Mueller2017Morphological}, but also utilize muscle
synergies~\cite{DAvella2003Combinations} and thus fewer control parameters. And while it is
difficult to clearly attribute the effect of the embodiment~\cite{Zahedi2013Quantifying}, for us,
this gives reason to believe that the simplicity, and thus efficiency, comes from a
low-dimensionality of the controller~\cite[Sec.~5]{Montufar2015CheapControl}.

That this is indeed true for certain examples has been shown both in theory and practice
before~\cite{Datar2025Universal,Montufar2015Geometry,Li2018Measuring}. Li et
al.~\cite{Li2018Measuring} call this lower dimension ``intrinsic dimension'', 
which is the smallest dimension of a parameter subspace in which the model
still achieves comparable metrics. They show, for example, that
for achieving $90\,\%$ of the performance of a fully connected network on MNIST,
the dimensionality can be reduced by $99.6\,\%$. For a convolutional neural
network, it is still a $99.3\,\%$ reduction.
Others show that the ``least complex model which is able to learn
an optimal policy does learn faster''~\cite[Sec.~6]{Montufar2015Geometry}. For a neural network, being
more parameter-efficient can be done in different ways, some of which we look at in
\cref{chap:inflation}. In comparison, for LQR, where the controller is a gain matrix, the
corresponding approach is reducing its rank, for example via a nuclear-norm regularization~\cite{Recht2010Guaranteed}.
However, while these approaches are more parameter-efficient at achieving similar or slightly worse performances, it
was also observed that they tend to perform even worse when encountering
noise~\cite{Datar2025Universal}.

Robustness to noise, however, is one of the important characteristics of controllers acting under
real-world conditions~\cite{Rao2001Naive,Roy2021MachineLearning}. One approach to achieving this
in nature is the principle of ``redundancy and diversity''~\cite[p.~828]{Kitano2004Biological}, which stands
in direct contrast to our goal of reducing redundant parameterizations. Another reason for producing
robust controllers is the difficult transfer from a digital training environment to real-world
embodiments, commonly referred to as the ``reality gap''~\cite{Tobin2017Domain,Peng2018SimToReal}.
Approaches to training robust controllers typically introduce noise already during training, be it by modelling
measurement and disturbance noise explicitly, as the linear quadratic Gaussian (LQG) controller
does~\cite[p.~89]{Rao2001Naive}, or by including variations to the
dynamics~\cite{Peng2018SimToReal}.
These approaches, however, do not look at low-dimensionality, which makes sense as
Kitano~\cite[p.~833]{Kitano2004Biological} proposes that robustness cannot be increased
without sacrificing performance or accepting ``increased resource
demands''.

Before taking a look at the more complex case of RL, we start simpler by
using the simpler and better-understood LQR. This approach has been taken
before~\cite{Recht2019Tour} and does make sense given the guarantees regarding convergence and
effectiveness of policy optimization for it~\cite{Fazel2018Global,Takakura2023Structured}, as well
as the understanding of its geometry~\cite{Talebi2024Policy}.
It also offers methods which produce certain structures in its controller, like sparsity or our desired
low-dimensionality, as given by structured policy
iteration (S-PI)~\cite{Park2020Structured}. Approaches can generally be categorized into two groups if the underlying dynamics
are unknown and only data is available: estimating the model and then applying a model-based solver (indirect), or using the data
to estimate the gain matrix (direct)~\cite{Zhao2024DataEnabled}. Both the indirect, as in system identification
with uncertainty bounds~\cite{Dean2020Sample}, and the direct setting, with an $\Hinf$ constraint
that bounds the closed-loop gain from disturbance to output~\cite{Zhang2021Policy},
have been studied with regard to robustness, but, to our knowledge, not in combination with low-rank policies.

Research evaluating RL on control tasks has also shown that there is room for simplification, for
example using linear policies instead of complex, fully connected neural
networks~\cite[Sec.~4.2]{Mania2018Simple}, \cite{Rajeswaran2017Towards}. This, together with the
fact that parameters in neural networks are highly dependent, and thus redundant~\cite{Denil2013Predicting}, allows us to assume that
low-dimensionality can improve efficiency here. Examples of this exist, for example by factorizing a discretized low-rank
policy into a thin and wide matrix~\cite{Rozada2024Matrix} or by introducing an orthogonal projection inside the layer
structure~\cite{Todorov2026Orthogonal}. As a result, the dimensionality of the policy is reduced. A different approach uses ``information bottlenecks'',
which exhibit better generalization behavior in unseen environments~\cite[Sec.~4.2]{Igl2019Generalization},
which again could help us with robustness.

A clear case where low-dimensional dynamics underlie a high-dimensional observational space is in perception-based tasks.
Agents observing only images instead of the states of simple control tasks are faced with the challenge of extracting
these states from a high-dimensional, often nonlinear encoding~\cite[Sec.~1]{Watter2015Embed}. Many approaches, based on
state representation learning (SRL), exist, which use different variants to create a latent embedding of the image,
used as a state~\cite{Lesort2018State}. The most interesting approach, especially for us looking at embodied agents, is
the usage of so-called ``priors'', which restrict the latent space by using knowledge about the
environment~\cite[Sec.~2.2]{Lesort2018State}.
This is similar to Ay's argument that the embodiment actively defines the action space~\cite[Sec.~1]{Ay2015Geometric}.
Robotic priors were introduced by Jonschkowski and Brock~\cite{Jonschkowski2015Learning} in 2015, and then further
refined by Botteghi et al.~\cite{Botteghi2021LowDimensional} for continuous action spaces. A more specific version
restricted to position and velocity also exists here, training a model to encode an image into a low-dimensional
position state, from which the velocity is calculated by finite differences~\cite{Jonschkowski2017PVE}.
However, perception-based control is also affected by the problems of robust control,
as noise induced through perception directly affects the performance~\cite[Sec.~3.3]{Xu2021Perception}
and therefore requires approaches for mitigating them~\cite{Dean2020Robust}.

Our goal for this thesis is therefore to look a bit deeper at the combination of low-dimensionality
and robustness in an empirical way, motivated by the findings of Datar and Ay~\cite{Datar2025Universal}. LQR
serves as the perfect, well-investigated basis~\cite{Recht2019Tour}, offering low-rank approaches that have been proposed
and analyzed theoretically, but have not yet been applied in practice and under noise~~\cite[Sec.~4]{Park2020Structured}. Previously
investigated policies for LQR under noise in comparison did not look at low-dimensionality
specifically~\cite{Zhang2021Policy}.
For RL, low-dimensionality has been explored
previously~\cite{Agarwal2020FLAMBE,Stojanovic2024ModelFree}, and also
in comparison with robustness~\cite{Li2025AdaRL}, but this research does not take into account the embodiment we are interested in.
Closing this gap, our main interest is to compare robustness properties for low-dimensional policies for embodied
agents under increasingly more complex and noisier environments.

% ------------------------------------------

\section{Experimental Setup}
\label{sec:setup}

We now introduce with \cref{subsec:modified-loop} the modifications we apply to the
double integrator loop for our experiments, followed by the metrics in \cref{subsec:robustness}
we evaluate it on. The section is concluded with \cref{subsec:protocol}, where we state
our evaluation procedure.

\subsection{The Modified Loop}
\label{subsec:modified-loop}

\subsubsection{An Inflated Model}
\label{subsubsec:inflated-model}

To allow us to investigate a controller that has many redundant inputs and outputs, acting on a
system with low-dimensional dynamics, we first need to create such a setup for our previously introduced
double integrator. We do this by keeping the dynamics, so the state progression, like the loop in \cref{subsec:double-integrator},
but projecting the observation into a higher dimension, as well as projecting actuation down from such.

More specifically, we show the modification to our loop with two random orthonormal matrices $U_1$
and $U_2$. For an inflated input dimension $d_{\mathrm{in}}$ and output dimension $d_{\mathrm{out}}$,
the two lifts are constructed the following way:
\begin{equation}
    \label{eq:lifts}
    U_1 \in \R^{d_{\mathrm{in}} \times n}, \quad U_1\T U_1 = I_n,
    \qquad
    U_2 \in \R^{m \times d_{\mathrm{out}}}, \quad U_2 U_2\T = I_m,
\end{equation}
where $U_1$ is the orthonormal factor of the QR decomposition of a random matrix, whose
entries are sampled from the standard normal distribution. $U_2$ is the transpose of such a factor
sampled from a separately drawn matrix. Both are drawn once per seed, and then kept fixed throughout.
A single step of the loop then becomes
\begin{subequations}
\label{eq:inflated-loop}
\begin{align}
    y_k     &= U_1 x_k , \label{eq:inflated-obs}\\
    v_k     &= -K y_k , \qquad K \in \R^{d_{\mathrm{out}} \times d_{\mathrm{in}}} , \label{eq:inflated-ctrl}\\
    u_k     &= U_2 v_k , \label{eq:inflated-act}\\
    x_{k+1} &= A x_k + B u_k = A x_k + B U_2\, v_k . \label{eq:inflated-dyn}
\end{align}
\end{subequations}
Compared with \eqref{eq:ss-cont}, the output matrix thus becomes $C = U_1$, so that $q = d_{\mathrm{in}}$,
and the controller acts through the input matrix $B U_2 \in \R^{n \times d_{\mathrm{out}}}$.

This way, the controller only sees the inflated observation and acts in a higher dimension. As any actuation
effort is penalized, its weight $\tilde{R}$ acts on the inflated actuation and therefore has the dimension
$d_{\mathrm{out}} \times d_{\mathrm{out}}$. To
simplify things, we also have the state weight $\tilde{Q}$ in this lifted dimension, acting on the observation. Here and in the
following, a tilde marks the inflated counterpart of a quantity of the uninflated problem. As $U_1$ has
orthonormal columns, and we use $\tilde{Q} = I_{d_{\mathrm{in}}}$ throughout, this simplifies to
\begin{equation}
    \label{eq:lifted-state-cost}
    \begin{aligned}
    \left( U_1 x_k \right)\T \tilde{Q} \left( U_1 x_k \right)
    &= x_k\T U_1\T U_1\, x_k \\
    &= x_k\T I_n\, x_k \\
    &= x_k\T Q\, x_k , \qquad Q = I_n ,
    \end{aligned}
\end{equation}
with the uninflated $Q$ of \cref{subsec:lqr}, so it indeed still penalizes how far our mass is from zero. The cost in the modified loop then
becomes
\begin{equation}
    \label{eq:inflated-cost}
    \begin{aligned}
    J(x_0, K)
    &= \sum_{k=0}^{\infty} \left( y_k\T \tilde{Q}\, y_k + v_k\T \tilde{R}\, v_k \right) \\
    &= \sum_{k=0}^{\infty} \left( x_k\T U_1\T \tilde{Q}\, U_1 x_k + x_k\T U_1\T K\T \tilde{R} K U_1 x_k \right) ,
    \end{aligned}
\end{equation}
with the closed loop $x_{k+1} = (A - B U_2 K U_1)\, x_k$.

Throughout the thesis, we use the following synonyms for different $d_{\mathrm{in}}$ and
$d_{\mathrm{out}}$ dimensionalities, applied to the loop, which can be found in \cref{tab:inflation-sizes}.
\begin{table}[htbp]
    \centering
    \caption{Preset inflation sizes.}
    \label{tab:inflation-sizes}
    \begin{tabular}{@{}lrr@{}}
        \toprule
        Name & $d_{\mathrm{in}}$ & $d_{\mathrm{out}}$ \\
        \midrule
        \emph{base}   & 2   & 1   \\
        \emph{small}  & 10  & 5   \\
        \emph{medium} & 50  & 25  \\
        \emph{large}  & 200 & 100 \\
        \emph{pixel}  & 18\,432 & 256 \\
        \bottomrule
    \end{tabular}
\end{table}

\subsubsection{Noise Injection}
\label{subsubsec:noise-injection}

Now, the second thing we want to take a look at is robustness to
small changes or variations in the environment.
We thus introduce four different types of disturbances to be included in our model,
commonly known as noise. They are distinguished by the point at which they
affect the loop: what the controller \emph{observes} (measurement noise),
what its output actually \emph{does} (control noise), what the plant
\emph{is} (parametric uncertainty), and what else \emph{acts} on it
(external disturbance). Each of them can be activated separately, so that we
can clearly isolate the cause of an effect. Throughout this section
$k$ indexes the sampling instants, and $\sigma$ denotes the standard deviation,
or amplitude of the respective noise channel.

\paragraph{Measurement noise.}
Here we influence the sensors used to measure
the model state, such that capturing the same state twice does
not lead to the same observation $y_k$.

Let $n^{\mathrm{m}}_k$ denote the measurement noise at step $k$. Then
\begin{equation}
    \label{eq:measurement-noise}
    n^{\mathrm{m}}_k \sim \mathcal{N}\!\left(0, \sigma_{\mathrm{m}}^{2} I_{d_{\mathrm{in}}}\right),
    \qquad
    y_k = U_1 x_k + n^{\mathrm{m}}_k ,
    \qquad \sigma_{\mathrm{m}} \in \R_{+} ,
\end{equation}
drawn independently at every step.

\paragraph{Control noise.}
As the name says, acts on the applied control input. With it,
we try to model the difference between the intended movement of an actuator
and its actual action due to different factors. Here we deal with
\emph{multiplicative} noise, which scales with the input. In
detail this means that the applied controller output $\tilde{v}_k$ is obtained as
\begin{equation}
    \label{eq:control-noise}
    n^{\mathrm{c}}_k \sim \mathcal{N}\!\left(0, \sigma_{\mathrm{c}}^{2} I_{d_{\mathrm{out}}}\right),
    \qquad
    \tilde{v}_k = \left( \mathbf{1}_{d_{\mathrm{out}}} + n^{\mathrm{c}}_k \right) \odot v_k ,
\end{equation}
where $\odot$ denotes the elementwise product, so that $\sigma_{\mathrm{c}}$ is
the \emph{relative} error. Applying this to the input, the next step becomes
\begin{equation}
    \label{eq:control-noise-step}
    x_{k+1} = A\, x_k + B U_2\, \tilde{v}_k .
\end{equation}

\paragraph{Parametric uncertainty.}
This noise acts directly on the matrices of the model. Here, we either decide to keep
it time invariant, meaning $\Delta^{A}_k = \Delta^{A}_0$ and $\Delta^{B}_k = \Delta^{B}_0$ for all $k$, or to resample it
at every step. The former models a wrongly identified system, while the latter
simulates one that continuously changes. As with measurement noise, it is additive.
Applied, it acts on the system the
following way:
\begin{equation}
    \label{eq:parametric-noise}
    \left( \Delta_k^{A} \right)_{ij},\ \left( \Delta_k^{B} \right)_{ij}
    \overset{\text{i.i.d.}}{\sim} \mathcal{N}\!\left(0, \sigma_{\mathrm{p}}^{2}\right),
    \qquad
    x_{k+1} = \left( A + \Delta_k^{A} \right) x_k + \left( B + \Delta_k^{B} \right) U_2\, \tilde{v}_k .
\end{equation}
Worth noting is that resampling every
step leaves the model correct \emph{on average}, since
$\E[\Delta^{A}_k] = \E[\Delta^{B}_k] = 0$, and mainly adds variance to the trajectory. This stands
in contrast to the time-invariant version, where we have a discrepancy between
what we compute actions on and what actually happens.

\paragraph{External disturbance.}
This noise models external influences, like a wind pushing on our
agent, or rain
hitting it with some frequency. It only affects the position, so with the first unit
vector $e_1 \in \R^{n}$ and $\sigma_{\mathrm{d}}$ as the noise amplitude, we get
\begin{equation}
    \label{eq:dist-sin}
    n^{\mathrm{d}}_k = \sigma_{\mathrm{d}} \sin(0.01\,k)\, e_1 .
\end{equation}
As the disturbance acts on the true state $x_k$, it is applied before any
lifting of the observation~\eqref{eq:inflated-obs}. The loop can thus
be described as
\begin{equation}
    \label{eq:disturbance-step}
    x_{k+1} = A\, x_k + B U_2\, \tilde{v}_k + n^{\mathrm{d}}_k .
\end{equation}

\subsubsection{The Full Loop}
\label{subsubsec:full-noise}

While we, until now, have only regarded each noise type on its own,
there is also the possibility to include them all at once. Combining all the channels
with the inflated loop, a single step becomes
\begin{subequations}
\label{eq:noisy-loop}
\begin{align}
    y_k         &= U_1 x_k + n^{\mathrm{m}}_k , \label{eq:noisy-obs}\\
    v_k         &= -K y_k , \label{eq:noisy-ctrl}\\
    \tilde{v}_k &= \left( \mathbf{1}_{d_{\mathrm{out}}} + n^{\mathrm{c}}_k \right) \odot v_k , \label{eq:noisy-act}\\
    x_{k+1}     &= \left( A + \Delta_k^{A} \right) x_k + \left( B + \Delta_k^{B} \right) U_2\, \tilde{v}_k + n^{\mathrm{d}}_k . \label{eq:noisy-dyn}
\end{align}
\end{subequations}

% ---------------------------

\subsection{Metrics and Robustness}
\label{subsec:robustness}

In this thesis we try to answer two main questions. First, how well
does a trained controller perform generally, as well as under influences
unknown to it during training? Secondly, what influence does its size
have on its performance, and what defines sizes? Thus we aim to introduce
the metrics here to quantify these questions and our results.

\subsubsection{Cost and Suboptimality}
\label{subsubsec:suboptimality}

By \cref{subsubsec:lqr-optimal} no gain achieves a
lower cost \eqref{eq:cost-of-gain} than $K^{*}$ on the clean loop,
and every controller we train is therefore in the best case only as good as it
unless noise is involved.
This serves as our baseline, and we evaluate the cost produced by our controllers
in relation to it.

While we can easily calculate the infinite-horizon cost of a gain
via the Lyapunov equation \eqref{eq:lyapunov}, we can only apply it
to a closed-loop system. However, when adding noise to it,
this is not the case anymore. Additionally, the neural
networks we train in later chapters do not admit a gain-matrix-based
policy $K$, which would not allow cross-chapter comparison.

So, we instead evaluate them over a fixed horizon via rollout. Hence for
a given controller, a noise condition $c$, which is a subset of channels and an
intensity (\cref{subsubsec:noise-strength}), and a horizon $T$, we
simulate starting from a randomly drawn $x_0 \sim \mathcal{D}$.
The resulting loop \eqref{eq:noisy-loop} under $c$ thus accumulates
\begin{equation}
    \label{eq:empirical-return}
    J_T^{(i,s)}(K; c)
    = \sum_{k=0}^{T-1}
      \left( \left( U_1 x_k \right)\T \tilde{Q} \left( U_1 x_k \right) + \tilde{v}_k\T \tilde{R}\, \tilde{v}_k \right) ,
\end{equation}
where $s$ indexes the seed the run was trained under and $i$ the episode
within it.

It is important to note where $\tilde{Q}$ and $\tilde{R}$ act. As $\tilde{Q}$ penalizes the difference
from the origin, we apply it on the inflated state \emph{before} any noise is
applied. This way we can evaluate the actual position of the mass, as shown
in \eqref{eq:lifted-state-cost}. For $\tilde{R}$, we handle it differently. Here,
every actuation energy adds to the cost, irrespective of its origin (so whether the
controller commanded it or it is noise).

We first define the set of pairs $(s,i)$ whose episode stayed finite under both the
trained controller and the baseline,
\begin{equation}
    \label{eq:finite-pairs}
    \mathcal{F}_{S} \coloneqq \left\{ (s,i) \;:\; J^{(i,s)}_T(K; c) < \infty ,\ J^{(i,s)}_T(K^{*}; c) < \infty \right\} .
\end{equation}
Now we can state our first main metric, which
is the relative excess of the cost produced by a
trained controller over the cost of the optimal controller $K^{*}$
under the same condition. It is taken \emph{per episode}, meaning
an episode $(s,i)$ of a trained controller and episode $(s,i)$ of the baseline
are compared here. Both start from the same $x_0$ and are affected by
the same noise. Overall this yields
\begin{equation}
    \label{eq:suboptimality}
    S^{(i,s)}(K; c)
        = \frac{J^{(i,s)}_T(K; c) - J^{(i,s)}_T(K^{*}; c)}
                {J^{(i,s)}_T(K^{*}; c)} ,
    \qquad
    \hat{S}(K; c)
        = \frac{1}{\lvert \mathcal{F}_{S} \rvert}
          \sum_{(s,i) \in \mathcal{F}_{S}} S^{(i,s)}(K; c) ,
\end{equation}
where $S = 0$ means that the trained controller performed as well
as our baseline and $S = 0.1$ that it did $10\,\%$ worse.

\subsubsection{Robustness to Noise}
\label{subsubsec:robustness-metric}

Before introducing our second main metric, we need to define 
\begin{equation}
    \label{eq:empirical-cost}
    \hat{J}_T(K; c)
    = \frac{1}{\lvert \mathcal{F} \rvert}
      \sum_{(s,i) \in \mathcal{F}} J^{(i,s)}_T(K; c) ,
    \qquad
    \mathcal{F} = \left\{ (s,i) \;:\; J^{(i,s)}_T(K; c) < \infty \right\} ,
\end{equation}
which is the cost averaged over all episodes. In the following we will also
use its variation $\hat{J}^{(s)}_T(K; c)$, where we pool over all episodes
of a seed $s$.

Our second metric relates to robustness, and helps us evaluate how much a controller is influenced
by a noise condition $c$. It is the ratio of the test cost under noise $c$ to the test cost under
the clean loop, which we from now on will call \emph{relative degradation}. Formalized, this gives
\begin{equation}
    \label{eq:relative-degradation}
    \kappa^{(s)}(K; c) \coloneqq \frac{\hat{J}^{(s)}_T(K; c)}{\hat{J}^{(s)}_T(K; \varnothing)} ,
\end{equation}
where $\hat{J}^{(s)}_T$ is \eqref{eq:empirical-cost} containing only the episodes of seed $s$.
Thus, the closer $\kappa^{(s)}(K; c)$ is to $1$, the more robust a controller is. 
We can also compare the relative degradation to the performance of the optimal gain, which we calibrated in \eqref{eq:noise-calibration}
to specific ratios, by taking the ratio of the costs averaged over all seeds,
\begin{equation}
    \label{eq:relative-degradation-pooled}
    \kappa(K; c) \coloneqq \frac{\hat{J}_T(K; c)}{\hat{J}_T(K; \varnothing)} ,
\end{equation}
for which $\kappa_{\ell} = \kappa(K^{*}; c_{\ell})$ holds by \eqref{eq:noise-calibration}.
As a result, this gives us the possibility to compare whether a specific trained
controller is more or less robust than the baseline, or any other controller.

However, this metric is misleading when comparing a controller which has performed
very badly on the clean loop with one that performed well.
If we assume that both achieve the same $\kappa(K; c)$, this would equate
to a much larger absolute increase in the test cost for the former, compared
to the latter. To quantify this, we introduce the following metric,
\begin{equation}
    \label{eq:absolute-degradation}
    \frac{\Delta \hat{J}}{\hat{J}^{*}}(K; c) \coloneqq \frac{\hat{J}_T(K; c) - \hat{J}_T(K; \varnothing)}{\hat{J}_T(K^{*}; \varnothing)} ,
\end{equation}
and call it \emph{absolute degradation}, indicating the increase in cost caused by noise,
in units of the optimal clean cost. For $K^{*}$, it equals $\kappa_{\ell} - 1$, i.e.\ $0$ on the
clean loop.

\subsubsection{The Size of a Policy}
\label{subsubsec:policy-size}

We have often mentioned ``low-dimensional'' controllers until now; however,
we want to make the number of parameters comparable, not only
between different LQR controllers, but also across LQR and RL.
So it is indeed a question of size. For a neural network we can
simply count the number of weights and biases of the actor. For a gain matrix, however,
we need to find a definition that is comparable to it.

The simplest approach for a $K \in \R^{d_{\mathrm{out}} \times d_{\mathrm{in}}}$ is the number of its entries,
$d_{\mathrm{out}} \cdot d_{\mathrm{in}}$. In \cref{chap:lqr}, however, we reduce the rank of $K$.
Optimizing the storage size for $K$ of rank $r$, the singular value decomposition~\cite[p.~85, Thm.~13.1]{Werner2021Optimal}
separates it into
\begin{equation}
    \label{eq:svd}
    K = \Phi \operatorname{diag}(\chi_1, \dots, \chi_r) \Psi\T ,
    \qquad
    \Phi \in \R^{d_{\mathrm{out}} \times r}, \quad
    \Psi \in \R^{d_{\mathrm{in}} \times r} ,
\end{equation}
with $\Phi$ and $\Psi$ having orthonormal columns and $\chi_1 \geq \dots \geq
\chi_r > 0$ being the non-zero singular values. Hence, we only need to
store $d_{\mathrm{out}} \cdot r$ entries for $\Phi$, $r$ for the singular values and
$r \cdot d_{\mathrm{in}}$ for $\Psi$, which in total yields $r(d_{\mathrm{out}} + d_{\mathrm{in}} + 1)$.
For a small enough $r$ this is less, so we take the minimum of both:

\begin{definition}[Policy size]
\label{def:policy-size}
The size of a linear policy $K \in \R^{d_{\mathrm{out}} \times d_{\mathrm{in}}}$ of rank
$r = \rank(K)$ is
\begin{equation*}
    \lvert K \rvert \coloneqq \min \left\{ d_{\mathrm{out}} \cdot d_{\mathrm{in}}, \; r \, (d_{\mathrm{out}} + d_{\mathrm{in}} + 1) \right\} ,
\end{equation*}
the number of real numbers required to store it.
The size of a neural policy is its total number of parameters in the
actor, counting both fixed and trainable ones.
\end{definition}

\subsubsection{Noise Intensity}
\label{subsubsec:noise-strength}

\Cref{subsubsec:noise-injection} introduced four channels which change the system
compared to the base model. Each noise channel is separated further
into five intensities,
\begin{center}
    \emph{none} $\prec$ \emph{weak} $\prec$ \emph{standard} $\prec$
    \emph{strong} $\prec$ \emph{divergent},
\end{center}
where \emph{none} leaves the system as it is and thus serves as our
control condition. All the other intensities are calibrated by how much
they increase the cost of the optimal gain $K^{*}$ on the uninflated system.
For each channel and intensity $\ell$, we choose the standard deviation
in a way such that for the condition $c_{\ell}$ under the test protocol of \cref{subsec:protocol}
\begin{equation}
    \label{eq:noise-calibration}
    \hat{J}_T(K^{*}; c_{\ell}) = \kappa_{\ell}\, \hat{J}_T(K^{*}; \varnothing),
    \qquad
    \kappa_{\ell} \in \{1.2,\ 1.5,\ 3,\ 10\}
\end{equation}
applies for \emph{weak}, \emph{standard}, \emph{strong} and \emph{divergent}
respectively.
For $K^{*}$, the increase stays the same under inflation for all channels
except control noise (\cref{tab:noise-size-ratios}). As it acts elementwise on the inflated output, its
effect on the plant is averaged over the number of actuators: \emph{divergent} control noise raises the
cost of $K^{*}$ by a factor of $9.97$ without inflation, but only by $6.80$ and
$6.62$ at the sizes $d_{\mathrm{in}} \times d_{\mathrm{out}} = 50 \times 25$ and $200 \times 100$.
For simplicity, however, we do not vary the levels by size, but instead accept this
as a side effect.
The concrete standard deviations per channel are listed in
\cref{tab:noise-levels}.

\subsection{Test Protocol}
\label{subsec:protocol}

Having introduced our metrics, we now define a protocol for how
we train and compare the different solver approaches with these metrics.

First, we note that we always use $\tilde{Q} = I_{d_{\mathrm{in}}}$ and $\tilde{R} = I_{d_{\mathrm{out}}}$ (so $Q = I_n$ and $R = I_m$ on the uninflated system), which weighs position,
velocity and control effort equally. Initial states are drawn uniformly from the
box $[-10, 10] \times [-1, 1]$. Some solvers require an initial stabilizing gain,
which for the uninflated system is $K_0 = [\,0.1,\ 0.5\,]$, whose closed loop has $\spr(A - B K_0) \approx 0.9975$.
On an inflated system, we use
\begin{equation}
    \label{eq:initial-gain}
    \tilde{K}_0 = U_2\T K_0 U_1\T + K_{\perp} , \qquad U_2 K_{\perp} U_1 = 0 , \qquad \lVert K_{\perp} \rVert_F = 0.5 .
\end{equation}
We create the random matrix $K_{\perp}$ such that $U_2 K_{\perp} U_1 = 0$, and thus adding it does not change
the stability property of the matrix. This way, the effective gain remains $U_2 \tilde{K}_0 U_1 = K_0$,
while $\tilde{K}_0$ is full rank, which offers the possibility of rank reduction.

For each algorithm, we train ten different controllers, each under a different fixed seed,
until convergence. We define convergence for LQR solvers as follows:
\begin{definition}[Convergence]
\label{def:convergence}
Let $g_j$ be the signal that a solver observes at iteration $j$, and let $\bar{g}_j$ be its
mean over the last $N_{\mathrm{w}}$ iterations. The solver has \emph{converged} at the first
iteration $j$ for which
\begin{equation*}
    \bar{g}_j \ge \left( 1 - \varepsilon_{\mathrm{rel}} \right) \bar{g}_{j - N_{\mathrm{p}}} ,
\end{equation*}
i.e.\ once the average has not decreased by the relative tolerance $\varepsilon_{\mathrm{rel}}$
within the last $N_{\mathrm{p}}$ iterations. After an iteration with an
infinite signal (caused by $K \notin \mathcal{K}$, \eqref{eq:stabilizing-set}), the comparison
restarts, such that we never compare with a window across an infinite iteration.
\end{definition}

The signal is something the solver itself observes and optimizes for,
like the cost \eqref{eq:cost-of-gain}, an estimate, or a regularized
version of it for solvers introducing structure. Some also have
additional early-stopping criteria. They will be introduced
in detail in \cref{chap:lqr}.

For all of them, we use the best iterate by
$\bar{g}_j$ to identify the policy we want to use.

The parameters of \cref{def:convergence} for the model-based and the model-free solvers are
listed in \cref{tab:convergence-lqr}.
\begin{table}[htbp]
    \centering
    \caption{Parameters of the convergence rule (\cref{def:convergence}) for the LQR solvers.}
    \label{tab:convergence-lqr}
    \begin{tabular}{@{}lrr@{}}
        \toprule
        Parameter & Model-based solvers & Model-free solvers \\
        \midrule
        Window $N_{\mathrm{w}}$                   & 1         & 50        \\
        Patience $N_{\mathrm{p}}$                 & 50        & 100       \\
        Relative tolerance $\varepsilon_{\mathrm{rel}}$ & $10^{-4}$ & $10^{-3}$ \\
        \bottomrule
    \end{tabular}
\end{table}

The number of steps the plant progressed and solver iterations until convergence are also
counted and reported. An upper cap of $20\,000$ iterations is set to prevent looping.

For RL-based algorithms we do not use such early-stopping logic. Following our experiments in
\cref{sec:app-rl-convergence}, we observed that implementing such a rule would lead to worse
policies in the end. Instead, the best policy is achieved close to the end, which leads us
to running each algorithm for as many steps as noted in \cref{tab:rl-step-caps}, and then
taking the final policy. Every $25\,000$ environment steps throughout, we
evaluate the current policy's performance, to track its training process.
\begin{table}[htbp]
    \centering
    \caption{Environment step budget of the RL solvers per inflation size.}
    \label{tab:rl-step-caps}
    \begin{tabular}{@{}lrrrr@{}}
        \toprule
        & \emph{small} & \emph{medium} & \emph{large} & \emph{pixel} \\
        \midrule
        Budget (environment steps) & $10^{7}$ & $1.5 \cdot 10^{7}$ & $3 \cdot 10^{7}$ & $10^{7}$\textsuperscript{$*$} \\
        \bottomrule
    \end{tabular}

    \smallskip
    {\footnotesize \textsuperscript{$*$}In the camera environment of \cref{sec:pixel-env}, one step corresponds to one decision
    of the agent, which is then held for $30$ environment steps. There, the step budget is $10^{7}$ agent steps, which equals
    $3 \cdot 10^{8}$ environment steps.\par}
\end{table}

After our controllers have been trained, we evaluate them per seed for \emph{100 episodes}
each, running for \emph{$T = 2000$ steps} per episode with a sampling period of $h = 0.01\,\mathrm{s}$.
For each episode we calculate the empirical cost \eqref{eq:empirical-return} and its
suboptimality \eqref{eq:suboptimality} in comparison to the optimal gain run with the same initial state
$x_0 \sim \mathcal{D}$ and sampled noise. In the end, for each solver architecture,
we get a distribution of suboptimalities over a total of 1000 trajectories.

The inflated size is always the same between training and evaluation. For noise
experiments, however, the controller is trained on a clean system, and then exposed to one with noise
activated during testing.

On the clean system, some trained controllers might not be able to stabilize the system,
and are therefore not functional. Those are therefore removed from the metric calculations
and instead reported separately. To define a diverging controller, we differentiate two cases:
\begin{enumerate}
    \item Linear gain: the controller becomes unstable, i.e.\
    $\spr(A - B U_2 K U_1) \geq 1$.
    \item Neural network: the controller in seed $s$ is not able to control the mass,
    i.e.\ contain it in the range of initial positions at the end of any of
    its episodes, formally once $\lvert p^{(i,s)}_T \rvert > 10$ for some episode $i$.
\end{enumerate}

Under the stronger noise conditions, individual episodes might still diverge, even
though their controller was stable before. To keep our mean finite, we exclude them
from the calculation, and report them separately as the divergence rate.
\begin{equation}
    \label{eq:divergence-rate}
    D(K; c) \coloneqq \frac{1}{\lvert \mathcal{S} \rvert\, N}
    \left\lvert \left\{ (s, i) \;:\; J^{(i,s)}_T(K; c) = \infty ,\ J^{(i,s)}_T(K^{*}; c) < \infty \right\} \right\rvert ,
\end{equation}
where $\mathcal{S}$ is the set of seeds whose controller is stable on the clean loop and $N = 100$ the
number of episodes per seed.

% ------------------------------------------

\section{Approach}
\label{sec:approach}

As mentioned before, the LQR serves as the perfect starting example,
which is exactly why we start our investigation here. For our investigations, we will
use the double integrator, as mentioned in \cref{subsec:double-integrator}.

First, as a basis, we will start by comparing
the standard policy gradient (PG) approach~\cite{Fazel2018Global} with S-PI~\cite{Park2020Structured}
on the true model (\emph{TrP} and \emph{TrS}). Their relative performance to the baseline will be evaluated under the different
inflation sizes introduced in \cref{tab:inflation-sizes}. Then, the same controllers will be exposed
to the different noise channels and intensities we introduced.

Afterwards, we will take a look at the cases where we do not have access to the real model. We compare
the relative performance of both direct and indirect approaches to our baseline, which is
the optimal controller calculated via the DARE \eqref{eq:dare} on the true model. For the indirect
approaches, we will use numerical algorithms for subspace state-space system identification
(N4SID)~\cite{VanOverschee1994N4SID} to estimate the underlying model from one
trajectory. A more detailed description, as well as the exact parameters selected,
will be provided in the sections discussing the setup. With this estimated model as
the basis, we employ three different approaches, called N4SID + DARE (\emph{IdD}), N4SID + Model-based PG (\emph{IdP}) and N4SID + S-PI (\emph{IdS}).
They will be complemented by Model-free Natural Policy Gradient (\emph{MfN})~\cite{Fazel2018Global} and model-free S-PI (\emph{MfS})~\cite{Park2020Structured},
which are direct solvers and estimate the gradient from episode rollouts.
All these will also be discussed in \cref{chap:lqr}.

In \cref{chap:inflation}, we will then first introduce RL-based approaches. With proximal policy optimization
(PPO)~\cite{Schulman2017Proximal} being one of
the most used architecture variants in robotics~\cite{Rudin2022Learning,Andrychowicz2020Learning,Kaufmann2023Champion}, we choose this as our main architecture,
and apply different variations of it. First, by using only one layer and dropping the activation function and the biases,
we force the trained controller to be linear, which allows us to extract
it as a gain matrix that we can closely compare with our LQR approaches (Linear PPO, \emph{Lin}). The fully connected, standard version is named
Fully-connected PPO (\emph{MLP}), and will also be evaluated. Additionally, we will introduce structures that force the observations through a low-dimensional latent
space. Here, we will reuse the N4SID system identifier from before to identify the output matrix $\hat{C}$, before applying a change of
basis to it, to achieve orthonormality in its columns. We project
the high-dimensional observation through $\hat{C}\T$ into a lower dimension, and then train PPO with this as the input. Therefore, the name for it is State encoder + PPO (\emph{StE}).
A similar approach is used for Scalable encoder + PPO (\emph{ScE}), which is based on Jonschkowski and Brock~\cite{Jonschkowski2015Learning} and
Botteghi et al.~\cite{Botteghi2021LowDimensional}, and uses
robotic priors to derive a low-dimensional latent state via a trainable encoder. Again, the construction following this encoder is PPO. Lastly, we
implement the orthonormal bottlenecks introduced by Todorov and Sabatelli~\cite{Todorov2026Orthogonal} (Bottleneck PPO, \emph{OrB}). It consists of sampling a fixed matrix
$U_{\mathrm{b}}$ with $d_{\mathrm{b}}$ orthonormal columns, which is added at the end of a trainable feature extractor, and used to project the output into a $d_{\mathrm{b}}$-dimensional space. Like with the other
encoders, we follow it with our PPO architecture. The same feature extractor, but without the projection matrix $U_{\mathrm{b}}$,
followed by PPO, gives us FC-Extractor + PPO (\emph{FcE}).
From now on, we refer to all PPO variants by their short names.

All these low-dimensional approaches, if visualized, exhibit a ``bow-tie'' shape, which in nature can actually
improve robustness~\cite[p.~829]{Kitano2004Biological}. So this is indeed interesting to investigate. However,
first, like in \cref{chap:lqr}, we start by comparing these different approaches across the different inflation
sizes, without any noise applied, to see how they generally behave. Then, we introduce noise, and compare
our approaches under \emph{large} inflation dimensions.

Taking our results from the experiments on the artificially inflated dimensions, in \cref{chap:perception}, we create a new setup, still using the double
integrator as our underlying structure. Here, we will observe the mass via a simulated camera, and place a number of actuators
evenly distributed among angles on a half-circle. This will give us an actual high-dimensional input, as well as a high-dimensional
output, that is in reality governed by low-dimensional dynamics.

By inputting images into our existing solvers, the dimensionality increases drastically. To be able to cleanly separate effects caused
by the nonlinearity of the image construction from architecture-related effects, we will first artificially inflate our inputs and outputs 
to exactly the sizes that will later be used by the pixelated environment. This evaluation will only be done on \emph{ScE},
\emph{MLP}, \emph{OrB} and \emph{FcE}, as the other approaches either do not scale well or are not relevant to the comparison.

Then, using the same architectural constructions for these solvers, but applied to our new environment,
we compare them with the additionally introduced position-velocity encoder~\cite{Jonschkowski2017PVE} followed by PPO (\emph{PVE}),
which uses similar robotic priors to \emph{ScE}, but focuses
more on position, velocity and acceleration. It learns positional encodings from the observations, and then calculates the velocity state based
on finite differences. Here, it is quite clear that we use knowledge of the limitations of the embodiment to construct more optimal controllers,
again in a bow-tie shape. Finally, we will introduce this environment to noise as well, and observe the effects of it.
\cleardoublepage

\chapter{Robustness of LQR Solvers in Different Dimensions}
\label{chap:lqr}

\section{A Known Model}
\label{sec:known-model}

We start by comparing two approaches that have knowledge of the underlying model,
and thus do not need to estimate gradients but can calculate them exactly. This allows
us to clearly separate the benefits of a low-dimensional controller from estimation
errors.

\subsection{Setting}
\label{subsec:known-model-setting}

% mb pg
Fazel et al.~\cite{Fazel2018Global} introduced global convergence guarantees for policy gradient methods based on the
true model in 2019. The paper describes a method for calculating the gradient of a policy from
the system and cost matrices, and updating the policy accordingly. As the paper uses the same $u_k = -K x_k$
convention as we do, we can implement the loop directly in the following way.

% Def mb_pg
Starting from an initial, stabilizing gain $K_0 \in \mathcal{K}$ and the second moment $\Sigma_0$ of the initial
state distribution, every iteration $j$ first solves the discrete Lyapunov equation \eqref{eq:lyapunov}
for $P_{K_j}$ using the implementation from SciPy~\needcite{scipy}. The state correlation matrix
$\Sigma_{K_j}$ then follows from solving a second Lyapunov equation,
\begin{equation}
    \label{eq:state-correlation}
    \Sigma_{K_j} = (A - B K_j) \Sigma_{K_j} (A - B K_j)\T + \Sigma_0 ,
\end{equation}
where, for initial states drawn uniformly from $[-10, 10] \times [-1, 1]$ as described in \cref{subsubsec:lqr-cost-of-gain},
$\Sigma_0 = \operatorname{diag}\left( 10^2/3, \, 1^2/3 \right)$. The policy gradient
\begin{equation}
    \label{eq:pg-gradient}
    \nabla f(K_j) = 2 \left( \left( R + B\T P_{K_j} B \right) K_j - B\T P_{K_j} A \right) \Sigma_{K_j}
\end{equation}
is then used for a gradient step on the gain matrix. \Cref{alg:mb-pg} summarizes the loop.
\begin{algorithm}[htbp]
    \caption{Model-based PG}
    \label{alg:mb-pg}
    \begin{algorithmic}[1]
        \REQUIRE stabilizing gain $K_0 \in \mathcal{K}$, second moment $\Sigma_0$, step size $\eta$
        \FOR{$j = 0, 1, 2, \dots$}
            \STATE solve \eqref{eq:lyapunov} for $P_{K_j}$
            \STATE solve \eqref{eq:state-correlation} for $\Sigma_{K_j}$
            \STATE compute the policy gradient $\nabla f(K_j)$ by \eqref{eq:pg-gradient}
            \STATE update the gain by $K_{j+1} = K_j - \eta \nabla f(K_j)$
        \ENDFOR
    \end{algorithmic}
\end{algorithm}

The fixed learning rate $\eta$ is a hyperparameter, which we set to $\eta = 10^{-5}$ for all inflation sizes, based on
experiments whose results can be found in \cref{sec:app-learning-rate-mb-pg}. This
policy gradient approach serves us as the basic, fully-parameterized implementation for this experiment.

% mb spi

Now we also want to implement a solver that produces a low-dimensional controller. As in our
case the controller is characterized by the gain matrix, and for matrices dimensionality is ``often
[\ldots] expressed by means of the rank''~\cite{Recht2010Guaranteed}, we aim at learning a low-rank
gain matrix. However, to make this comparable later to RL approaches that rely on the number of parameters,
we measure the complexity of a gain by its policy size $\lvert K \rvert$ from \cref{def:policy-size}, which uses
the more efficient version between the full and the factored parameterization.

Recht et al.~\cite{Recht2010Guaranteed} have already shown that nuclear norm minimization introduces
low-rank solutions for convex matrix problems. However, as we are dealing with a non-convex optimization
problem~\cite{Fazel2018Global}, we cannot transfer the results over directly. Here, Park et al.~\cite{Park2020Structured}
offer us an approach to impose different structures on our policy during learning via regularizers
based on norms, one of them being the nuclear norm. We implement the model-based approach introduced there
with small adaptations in the following way.

% Def mb_spi
We reuse the computation of the gradient $\nabla f(K_j)$ from \cref{alg:mb-pg}.
For inducing the structure, we use an adaptive learning rate combined with a line search criterion.
Starting from the initial step size $\eta = 1 / (2 \lambda)$, each iteration takes a gradient step and applies the
proximal operator of the nuclear norm, which soft-thresholds the singular values,
\begin{equation}
    \label{eq:spi-prox}
    K^{+} = U \operatorname{diag} \left( \max\{\sigma_1 - \eta \lambda, 0\}, \dots,
    \max\{\sigma_{\bar{r}} - \eta \lambda, 0\} \right) V\T ,
    \qquad \bar{r} \coloneqq \min\{d_{\mathrm{out}}, d_{\mathrm{in}}\} ,
\end{equation}
where $U \operatorname{diag}(\sigma_1, \dots, \sigma_{\bar{r}}) V\T$ is the singular value decomposition of
the next policy candidate $K^{+} = K_j - \eta \nabla f(K_j)$. The candidate $K^{+}$ is accepted if it is stabilizing, $\spr(A - B K^{+}) < 1$, and
satisfies the sufficient decrease condition
\begin{equation}
    \label{eq:spi-sufficient-decrease}
    f(K^{+}) \leq f(K_j) - \eta \tr \left( \nabla f(K_j)\T G \right) + \frac{\eta}{2} \lVert G \rVert_F^2 ,
    \qquad G \coloneqq \frac{1}{\eta} \left( K_j - K^{+} \right) .
\end{equation}
Otherwise, the step size is reduced to $\beta \eta$ and the proximal step is repeated. The learning loop further can be
stopped earlier if the criterion $\lVert K_{j+1} - K_j \rVert_F \leq \varepsilon$ is satisfied.
\Cref{alg:mb-spi} summarizes the loop.
\begin{algorithm}[htbp]
    \caption{S-PI}
    \label{alg:mb-spi}
    \begin{algorithmic}[1]
        \REQUIRE stabilizing gain $K_0 \in \mathcal{K}$, second moment $\Sigma_0$, weight $\lambda$, factor $\beta$,
        tolerance $\varepsilon$
        \FOR{$j = 0, 1, 2, \dots$}
            \STATE compute $\nabla f(K_j)$ as in \cref{alg:mb-pg}
            \STATE set the initial step size $\eta = 1 / (2 \lambda)$
            \FOR{at most $20$ line search steps}
                \STATE compute the candidate $K^{+}$ by \eqref{eq:spi-prox}
                \IF{$\spr(A - B K^{+}) < 1$ and \eqref{eq:spi-sufficient-decrease} holds}
                    \STATE \textbf{break}
                \ENDIF
                \STATE reduce the step size, $\eta \leftarrow \beta \eta$
            \ENDFOR
            \STATE update the gain, $K_{j+1} = K^{+}$
            \IF{$\lVert K_{j+1} - K_j \rVert_F \leq \varepsilon$}
                \STATE \textbf{stop}
            \ENDIF
        \ENDFOR
    \end{algorithmic}
\end{algorithm}

As one can see, there are a lot of parameters that can be modified here. We set $\varepsilon = 10^{-8}$,
and use the same $\beta = 0.5$ as in the experiments section of Park et al.~\cite{Park2020Structured}.
For $\lambda$, as it defines the learning rate as well, we have run an experiment with various values,
which can be found in \cref{sec:app-spi-lambda}. A value of $\lambda = 5$ has shown the most promising results across the sizes.

One thing that we implemented differently than the paper is the cap on the number of line search iterations in line four, which we set to 20.
Some of our previous experiments had shown that if we leave it uncapped, for unstable updates which are impossible
to stabilize, the loop runs until the update produces no effect at all.

Now that we have defined our controllers and how they learn, we test them in comparison to the DARE-calculated
solution \eqref{eq:dare} from \cref{subsubsec:lqr-optimal}.

\subsection{Results}
\label{subsec:known-model-results}

We apply our solvers to the true models base and inflated versions, for now noise-free, as described in \needcite{cref inflation loop}. 
Specifically, as in \needcite{cref to test protocol}, we train 10 controllers each per solver, before evaluating them on the test episodes.
The results are then pooled.

The most interesting finding is, how much faster SPI converges compared to the Policy Gradient. While it achieves its best training
suboptimality within around 50 iterations, MB PG still falls drastically short of it after 450 iterations. We can also see in 
\needcite{table showing rank / complexity} that the reduction in complexity for all inflation sizes has worked well with SPI. 
For our test suboptimality, SPI also outperforms MB PG here by at least one order of magnitude. \needcite{across sizes suboptimality}
visualizes this.

% the table and figure from above

Overall, this looks very promising for SPI, still, our goal is to find out whether it also is a robust controller.

\subsection{Noise Experiments}
\label{sec:true-model-noise}

Now, we reuse the controllers we trained in \needcite{cref to subsec above} and evaluate them under the test protocol
with noise\needcite{cref to noise introduction} injected. 

When looking at the results, one can notice that our baseline controller, which was optimal for the clean loop, mostly
performs worse than our trained controllers, except for external disturbance. However, when only comparing MB PG with
SPI, we can see differences in their performance, depending on the type of noise that is applied, as well as whether the
loop is inflated or not. For parametric noise for example, MB PG performs on average better than SPI here, however also
has a higher variance. When looking at external disturbance under large inflation, one can notice that with increasing
noise it becomes clearer that SPI is more robust here compared to MB PG. While for weak and standard their total spread
still shows some overlap, for strong and divergent SPI shows a cleanly seperated suboptimality below MB PG.

% figure of comparison across the sizes for external disturbance

For control noise in comparison, we can see that here MB PG is more robust on average across settings. However, its worst case
performance here is worse than SPI, so it has an overall wider spread of suboptimalities. Still, both mostly outperform the
baseline controller. Measurement noise becomes an interesting case here. While both controllers outperform the baseline, under
no inflation, MB PG seems to be much more robust here. However, when looking at the controllers under large inflation with the
same noise level applied, suddenly SPI shows more robustness here with little variation. MB PG on the other hand is only on
average better than the baseline, but in some cases worse. \needcite{figure below} visualizes this.

% left base noise measurement divergent, right large noise measurement divergent

Overall, for the strong level and large inflation, this is how both controllers compare to the baseline and each other across the
noise channels \needcite{table below}.

% table true model noise large channels

\subsection{Discussion}
\label{sec:true-model-discussion}

Overall, we note that trained controllers are mostly at an advantage compared to statically calculated controllers in regards to
robustness. When comparing now trained fully parameterized and structured approaches, differences definitely appear. Changes in
the underlying model, as in Parametric Uncertainty, affect all controllers mostly the same, even though a fully parameterized
version might offer better robustness on average. Looking at the other cases, we find that low-dimensional controllers actually
offer better robustness in noise related to the inputs of the controller, especially the larger the input dimension. Measurement
noise shows this clearly, as the performance of a overparameterized controller degrades here with increasing size, while a
regularized controller stays mostly the same, or even shows improvement. However, related to the output dimensions, under control
noise reduced controllers perform worse on average, however still better in the worst case than the full controller. And under large
size the reduction in needed parameters is of two magnitudes of order, which shows the efficiency of the reduced controller.

\section{Unknown Model across Dimensions}
\label{sec:unknown-model}

Now we take a look at the more realistic case, where we do not know the underlying model, and compare different approaches.
Before we can get into the results, we first take a deeper look at the solvers that we use and how the experiments are
set up.

\subsection{Setting}
\label{subsec:unknown-model-setting}

First, we introduce our basis for the indirect model-free approach, which is based on system identification. N4SID,
introduced by Van Overschee and De Moor~\cite{VanOverschee1994N4SID}, uses ``the projection of the future outputs [of the system]
onto the past and future inputs and the past outputs''~\cite[Sec.~3]{VanOverschee1994N4SID}. These projections are then called
``a function of the system matrices and the input-output block Hankel matrices''~\cite[Sec.~3]{VanOverschee1994N4SID}, which
allows us to reconstruct exactly that from the results. Important are further the assumptions that the systems are
observable (\cref{def:observability}) and controllable (\cref{def:controllability}). For simplicity, we rely on the
implementation of N4SID from SIPPY~\needcite{TODO}, which we integrate into our own setup. This implementation allows us
to recover the system matrices $A$, $B$ and $C$, whereas we implemented the recovery of the cost matrices $Q$ and $R$
from scratch based on~\cite{VanOverschee1994N4SID}.

The collection of observation and action data runs the plant for $T$ steps under the initial gain, with a randomly
sampled excitation added on top of it,
\begin{equation}
    \label{eq:identification-input}
    v_k = -K_0 y_k + \sigma_{\mathrm{e}} w_k , \qquad w_k \sim \mathcal{N}(0, I_{d_{\mathrm{out}}}) ,
    \qquad k = 0, \dots, T - 1 .
\end{equation}

As mentioned in the paper, ``as $i$ grows larger, the approximation error grows smaller''~\cite[Sec.~5.3]{VanOverschee1994N4SID},
with $i$ being dependent on the length of our trajectory described above used for identification. The other variable $j$
defines the dimensions of the Hankel matrices, and thus how far the projections reach into the future and the past.
As we want to be computationally efficient while still achieving the best performance, we evaluated different values
for both, to find the optimal setting to use in our experiment. The full results of the selection can be found in the
\cref{sec:app-identification}. Specifically, we use the settings in
\cref{tab:identification-settings}.
\begin{table}[htbp]
    \centering
    \caption{Identification settings per inflation size.}
    \label{tab:identification-settings}
    \begin{tabular}{@{}lrr@{}}
        \toprule
        Name & Trajectory length $T$ & Hankel horizon \\
        \midrule
        \emph{base}   & 500  & 5 \\
        \emph{small}  & 500  & 5 \\
        \emph{medium} & 500  & 5 \\
        \emph{large}  & 8000 & 5 \\
        \bottomrule
    \end{tabular}
\end{table}

The resulting identified system then serves as a surrogate for the true model
for the model-based approaches from \cref{subsec:known-model-setting}. This thus produces results without
ever having access to the real model.

An attentive reader might now ask why we would even want to use these iterative approaches, if we could just
solve the DARE \eqref{eq:dare} directly on this estimation. And indeed we include it in our following comparison.
However, to make this work with our construction of inflated state models, we need some modification to N4SID.
This stems from the fact that we need to calculate $K = \left( R + B\T P B \right)^{-1} B\T P A$ from
\eqref{eq:optimal-gain}, hence $R + B\T P B$ needs to be invertible.
But for our inflated space, the estimated $B$ and $R$ have the same null space, making the matrix singular on that subspace.
So to be able to apply this solver, we eliminate the additional singular values of $B$ that exceed the estimated order
of the model, thus reducing the model in dimension. This way we do not access the model, but create encoding and
decoding matrices that map between our dimensionally reduced model and the actual state observation, as well as the
actuation output. For the exact implementation, please refer to the repository~\needcite{our repository}.

Additionally, we also use direct model-free approaches, which are variations of the model-based approaches
from \cref{sec:known-model}. Unlike before, they do not calculate the gradient directly, but estimate
it based on the returns of $N_{\mathrm{traj}}$ slightly perturbed versions of the current policy, simulated for $l$ steps. As one might
assume, more ``data'', could conversely to better estimation, and thus better policy updates. So, to
settle this, we evaluated different combinations of them in \cref{sec:app-sample-size}. The best
trade-off between computational efficiency and quality gave $N_{\mathrm{traj}} = 100$ and $l = 800$.

% mf npg
We use the model-free natural policy gradient (NPG) estimation described by Fazel et al.~\cite{Fazel2018Global},
which at its core consists of the loop in \cref{alg:mf-gradient}. Instead of the proposed single-point estimator,
we use a two-point estimator inspired by Duchi et al. and Mania et al.~\cite{Duchi2015Optimal,Mania2018Simple},
as this drastically reduced the variance of the estimates, and thus the instability of the controller.

\begin{algorithm}[htbp]
    \caption{Two-point gradient estimation}
    \label{alg:mf-gradient}
    \begin{algorithmic}[1]
        \REQUIRE gain $K \in \R^{d_{\mathrm{out}} \times d_{\mathrm{in}}}$, number of perturbation pairs $N_{\mathrm{traj}}$,
        rollout length $l$, smoothing radius $r$
        \FOR{$i = 1, \dots, N_{\mathrm{traj}}$}
            \STATE sample $\Xi_i$ uniformly from the matrices with $\lVert \Xi_i \rVert_F = r$
            \STATE sample an initial state $x_0 \sim \mathcal{D}$
            \STATE simulate $K + \Xi_i$ and $K - \Xi_i$ for $l$ steps from $x_0$, with stage costs $c_k^{\pm}$ and
            observations $y_k^{\pm}$
            \STATE $\hat{J}_i^{\pm} = \sum_{k=0}^{l-1} c_k^{\pm}$, \quad
            $\hat{\Sigma}_i^{\pm} = \sum_{k=0}^{l-1} y_k^{\pm} \left( y_k^{\pm} \right)\T$, \quad
            $\delta_i = \hat{J}_i^{+} - \hat{J}_i^{-}$
        \ENDFOR
        \STATE $\bar{\delta} = \frac{1}{N_{\mathrm{traj}}} \sum_{i=1}^{N_{\mathrm{traj}}} \delta_i$
        \STATE $\widehat{\nabla f}(K) = \frac{1}{N_{\mathrm{traj}}} \sum_{i=1}^{N_{\mathrm{traj}}}
        \frac{d_{\mathrm{out}} d_{\mathrm{in}}}{2 r^2} \left( \delta_i - \bar{\delta} \right) \Xi_i$
        \STATE $\hat{\Sigma}_K = \frac{1}{N_{\mathrm{traj}}} \sum_{i=1}^{N_{\mathrm{traj}}} \left( \hat{\Sigma}_i^{+} + \hat{\Sigma}_i^{-} \right)$
        \STATE $\hat{f}(K) = \frac{1}{N_{\mathrm{traj}}} \sum_{i=1}^{N_{\mathrm{traj}}} \frac{1}{2} \left( \hat{J}_i^{+} + \hat{J}_i^{-} \right)$
        \RETURN $\widehat{\nabla f}(K)$, $\hat{\Sigma}_K$, $\hat{f}(K)$
    \end{algorithmic}
\end{algorithm}

We further aim to clarify a couple of things about the given loop.
Subtracting the mean difference $\bar{\delta}$ does not change the expected value of the estimate, since
$\E[\Xi_i] = 0$, but lowers its variance. A rollout whose observation exceeds $\lVert y_k \rVert = 10^{5}$,
we count as diverged. It is therefore ended early and charged a penalty cost of $10^{10}$. 
Additionally, for inflated dimensions, $\hat{\Sigma}_K$ is singular, so we precondition with its
truncated pseudoinverse $\hat{\Sigma}_K^{\dagger}$ instead, which discards singular values below $10^{-6}$ times the largest one.

Our full learning algorithm is then given in \cref{alg:mf-npg}.
\begin{algorithm}[htbp]
    \caption{Model-free NPG}
    \label{alg:mf-npg}
    \begin{algorithmic}[1]
        \REQUIRE stabilizing gain $K_0 \in \mathcal{K}$, $N_{\mathrm{traj}}$, $l$, $r$, step size $\eta$
        \FOR{$j = 0, 1, 2, \dots$}
            \STATE compute $\widehat{\nabla f}(K_j)$ and $\hat{\Sigma}_{K_j}$ by \cref{alg:mf-gradient}
            \STATE update the gain by $K_{j+1} = K_j - \eta \, \widehat{\nabla f}(K_j) \, \hat{\Sigma}_{K_j}^{\dagger}$
        \ENDFOR
    \end{algorithmic}
\end{algorithm}

After some evaluations, we have decided on a smoothing parameter of $r = 0.5$, and a learning rate of $\eta = 0.01$, based
on the results in \cref{sec:app-learning-rate-mf-npg}.

% spi mf
Next, we implement a model-free version of S-PI that relies on the same gradient
estimation (\cref{alg:mf-gradient}) as Model-free NPG. Another crucial difference,
apart from the estimation of the gradient compared to the calculation, is the ``usage of a fixed stepsize
$\eta$, rather than an adaptive stepsize from a backtracking linesearch''~\cite[Sec.~3.1]{Park2020Structured}, as that
would require access to the function value $f(K) = \tr(P_K \Sigma_0)$. Before we look at how we select this step size
and $\lambda$, we define our learning loop in \cref{alg:mf-spi}.
\begin{algorithm}[htbp]
    \caption{S-PI (model-free)}
    \label{alg:mf-spi}
    \begin{algorithmic}[1]
        \REQUIRE stabilizing gain $K_0 \in \mathcal{K}$, $N_{\mathrm{traj}}$, $l$, $r$, fixed step size $\eta$, weight $\lambda$,
        tolerance $\varepsilon$
        \FOR{$j = 0, 1, 2, \dots$}
            \STATE compute $\widehat{\nabla f}(K_j)$ by \cref{alg:mf-gradient}
            \STATE compute $K_{j+1}$ by \eqref{eq:spi-prox}, with $\widehat{\nabla f}(K_j)$ in place of $\nabla f(K_j)$
            \IF{$\lVert K_{j+1} - K_j \rVert_F \leq \varepsilon$}
                \STATE \textbf{stop}
            \ENDIF
        \ENDFOR
    \end{algorithmic}
\end{algorithm}

Now to the selection of the parameters. Park et al.\ already note the difficulty of finding a fixed step size
that does not produce divergence of the method~\cite{Park2020Structured}. Also, the example shown only
uses the model-based approach and the regularizer to induce sparse structure, not the model-free variant combined
with low dimensionality. So during our first tests, we often ran into either unstable controllers, or ones that
did not induce any structure at all, hinting at a too low $\lambda$. If we remember the proximal step
\eqref{eq:spi-prox}, we can see that it depends on both the step size and $\lambda$, as every singular value is
mapped to
\begin{equation}
    \label{eq:spi-threshold}
    \sigma_i \mapsto \max\{ \sigma_i - \eta \lambda, 0 \} .
\end{equation}
Even after some improvement by using the two-point estimator (\cref{alg:mf-gradient}), it remained quite a difficult
topic. So rather extensive experiments on the variation of $\lambda$ and the step size can be found in \cref{sec:app-spi-mf}.
Based on those, $\lambda = 4000$ and a step size of $\eta = 3.5 \times 10^{-6}$ were selected.

\subsection{Results}
\label{subsec:unknown-model-results}

\subsection{Noise Experiments on Large Inflation}
\label{subsec:noise-large-inflation}

\subsection{Discussion}
\label{sec:id-model-discussion}

\section{Discussion}
\label{sec:ch4-discussion}

\cleardoublepage

\chapter{Robustness of Reinforcement Learning Approaches}
\label{chap:inflation}

Having seen how low-dimensionality affects robustness with a simple, LQR-based controller, we
now want to take a look at a setting which allows us to judge architectures used in more
complex tasks as well. Still, we will stay with our introduced double integrator system,
but change the controllers for now.

\section{PPO and Its Variants}
\label{sec:ppo-variants}

\subsection{Proximal Policy Optimization}
\label{subsec:ppo}

Before getting into our implementation details, we want to explain the basic construction
of PPO. Introduced by Schulman et al.~\cite{Schulman2017Proximal} in 2017, it is an actor-critic algorithm,
which means it consists of an actor, guided by a policy, that is simulated in an environment,
and a critic, which tries to estimate the underlying value function based on the current policy
and the actor's returns~\cite{Konda2000ActorCritic}. In the case of PPO, as we are trying to optimize a policy
via policy gradient optimization, we ``[construct] an objective function whose gradient is the policy
gradient estimator''~\cite[Sec.~2.1]{Schulman2017Proximal}. What makes PPO special is the construction of this surrogate
function, based on ideas from Trust Region Policy Optimization~\cite{Schulman2015Trust}.
In the paper, two approaches are proposed: a Clipped Surrogate Objective or an Adaptive Kullback--Leibler (KL)
Penalty Coefficient. The former clips the probability ratio past a limit set by a hyperparameter,
imposing constraints on the policy update size~\cite[Sec.~3]{Schulman2017Proximal}, while the latter, rather than setting
a fixed KL for an update, sets a target KL divergence~\cite[Sec.~4]{Schulman2017Proximal}. In case further interest
arises, we refer to the original paper~\cite{Schulman2017Proximal}.
With the probability ratio $r_k(\theta)$ between the current policy $\pi_\theta$ and the policy
$\pi_{\theta_{\mathrm{old}}}$ that collected the data, the clipped objective becomes
\begin{equation}
    \label{eq:ppo-clip}
    \begin{aligned}
        r_k(\theta) &= \frac{\pi_\theta(v_k \mid y_k)}{\pi_{\theta_{\mathrm{old}}}(v_k \mid y_k)} , \\
        L^{\mathrm{CLIP}}(\theta) &= \hat{\E}_k \left[ \min\left\{ r_k(\theta) \hat{A}_k ,\;
            \operatorname{clip}\left( r_k(\theta), 1 - \varepsilon_{\mathrm{clip}}, 1 + \varepsilon_{\mathrm{clip}} \right) \hat{A}_k \right\} \right] .
    \end{aligned}
\end{equation}

\subsection{Training Setup}
\label{subsec:rl-training}

However, as has been noted before~\cite{Engstrom2020Implementation,Hsu2020Revisiting}, the actual implementation of PPO matters,
and can influence the effectiveness of it in practice. For our construction, we use
the architecture provided by Stable-Baselines3 (SB3)~\cite{StableBaselines2026PPO}.
As the standard architecture, for our Fully-connected PPO, we use the standard two-layer multilayer perceptron (MLP) with 64 neurons for both
the actor and the critic~\cite{StableBaselines2026Policy}.
Additionally, unless named otherwise, we will use the standard parameters from the SB3 library.
For the above named parts, we leave the clip range at the standard of $0.2$ and set the target KL to $0.02$.
Experimenting with different runtime lengths, we settled on the following values, which differ by inflation size:

% TODO: insert table here after experiments

Across all, the evaluation frequency, so when the policy of the actor is updated, is set to $25\,000$.

After initial runs with this basic setup, we could observe some divergence behavior, which
led us to consider additional implementation improvements, as suggested by Huang et al.~\cite{Huang2022Implementation}.
Specifically, we started using a vectorized environment, similar to the paper~\cite{Schulman2017Proximal}. Additionally,
we modified the discount factor to $\gamma = 0.999$ instead of $0.99$. However, the two most effective changes
were:
\begin{enumerate}
    \item Using learning rate annealing instead of a fixed rate~\cite[Detail~4]{Huang2022Implementation}. Now the learning
    rate is decreased with every update, such that at the end of the number of steps (\cref{tab:rl-step-caps}) it is equal
    to $0$. So this is
    \begin{equation}
        \label{eq:lr-annealing}
        \eta(N_{\mathrm{env}}) = \eta_0 \max\left\{ 1 - \frac{N_{\mathrm{env}}}{N_{\max}}, \; 0 \right\} .
    \end{equation}
    \item Separating the feature extractors for actor and critic~\cite[Detail~13]{Huang2022Implementation}. As described
    by~\cite{StableBaselines2026Policy}, to reduce computation, in the standard implementation,
    they share the extractor, which is ``a simple MLP network consisting of two layers of 64 neurons and Hyperbolic Tangent as the
    activation function''~\cite{Huang2022Implementation}.
    Later, when introducing pretrained extractors, they will be shared again, but for still learning ones, we will separate them.
\end{enumerate}
\Cref{alg:ppo} summarizes the resulting training loop.
\begin{algorithm}[htbp]
    \caption{PPO implemented here}
    \label{alg:ppo}
    \begin{algorithmic}[1]
        \REQUIRE initial actor $\pi_{\theta}$ and critic $V_{\phi}$, initial learning rate $\eta_0$, step cap $N_{\max}$,
        convergence parameters $N_{\mathrm{w}}$, $N_{\mathrm{p}}$, $\varepsilon_{\mathrm{rel}}$
        \STATE $N_{\mathrm{env}} \leftarrow 0$
        \WHILE{$N_{\mathrm{env}} < N_{\max}$}
            \STATE run $\pi_\theta$ for $2048$ steps, sampling $v_k \sim \pi_\theta(\cdot \mid y_k)$; store $y_k$, $v_k$ and the rewards
            \STATE normalize the rewards by a running estimate of the return's standard deviation
            \STATE compute the advantages $\hat{A}_k$ and value targets with $V_\phi$
            \STATE $\theta_{\mathrm{old}} \leftarrow \theta$, \quad set the learning rate $\eta(N_{\mathrm{env}})$ by \eqref{eq:lr-annealing}
            \FOR{epoch $= 1, \dots, 10$}
                \FOR{each minibatch of $64$ stored steps}
                    \STATE take an Adam step on $-L^{\mathrm{CLIP}}(\theta) + 0.5 \cdot (\text{value loss of } V_\phi)$,
                    with the gradient norm clipped to $0.5$
                \ENDFOR
                \IF{the approximate KL divergence to $\pi_{\theta_{\mathrm{old}}}$ exceeds $1.5 \cdot 0.02$}
                    \STATE \textbf{break}
                \ENDIF
            \ENDFOR
            \STATE $N_{\mathrm{env}} \leftarrow N_{\mathrm{env}} + 2048$
            \IF{$25\,000$ steps have passed since the last check}
                \STATE evaluate $\pi_\theta$, update $g_j$ with the mean training cost since the last check
                \IF{\cref{def:convergence} holds}
                    \STATE \textbf{stop}
                \ENDIF
            \ENDIF
        \ENDWHILE
        \STATE \textbf{return} the policy with the best smoothed signal $\bar{g}_j$
    \end{algorithmic}
\end{algorithm}

\subsection{Policy Architectures}
\label{subsec:rl-architectures}

Now, after having introduced the base construction, it makes it easier for us to describe variations of the same. We start with
Linear PPO, which allows us to compare the performance of PPO to the previous chapter. Here, our actor becomes a single layer,
with the bias removed, such that we can rewrite this network as
\begin{equation}
    \label{eq:linear-ppo}
    \mu(y_k) = W y_k = -K y_k , \qquad K \coloneqq -W \in \R^{d_{\mathrm{out}} \times d_{\mathrm{in}}} ,
\end{equation}
which allows us to extract a gain matrix $K$ that is comparable to the previous chapter. A second, already known, component is
used for State encoder + PPO. Here, we reuse an identified system from N4SID (\cref{sec:unknown-model}) to encode an observation into
a latent state:
\begin{equation}
    \label{eq:state-encoder}
    \hat{x}_k = \hat{C}\T y_k , \qquad \hat{C} \in \R^{d_{\mathrm{in}} \times \hat{n}} , \quad \hat{C}\T \hat{C} = I_{\hat{n}} .
\end{equation}
The input dimension for both the actor and the critic thus becomes $\hat{n}$, which is the order of the identified model.

However, as we will later scale our inputs to larger dimensions (\cref{chap:perception}), N4SID will become infeasible for such dimensions.
Still, we want to stay in line with the idea of encoding our observations into a latent state, as this exactly introduces the bow-tie structure
we want. So, we introduce another encoder, which offers this scalability, and call the combination Scalable encoder + PPO. This feature extractor
is based on the results from~\cite{Jonschkowski2015Learning,Botteghi2021LowDimensional}, which use robotic priors. Robotic priors, simply described,
are prior knowledge, like physical laws, from which we define loss functions. By adding them to the training process, and optimizing for them, we can
ensure that the learned representation is in line with these laws. As mentioned in \cref{sec:approach}, this allows us to make use of the embodiment
to improve our controller architecture. Here, we will use the following priors~\cite{Botteghi2021LowDimensional}:
\begin{itemize}
    \item simplicity: the state is lower in dimensionality than the observation,
    \item temporal coherence: state changes are rather slow,
    \item proportionality: same actions should show similar output strength,
    \item repeatability: same actions should show similar output changes,
    \item causality: TODO
\end{itemize}

We modify the loss calculations taken from the paper~\cite{Botteghi2021LowDimensional} slightly,
by dropping the Lreg term, and thus get the total loss function of
% TODO - cross check with code again
\begin{equation}
    \label{eq:prior-loss}
    \mathcal{L} = \omega_1 \mathcal{L}_{\mathrm{temp}} + \omega_2 \mathcal{L}_{\mathrm{prop}} + \omega_3 \mathcal{L}_{\mathrm{rep}}
        + \omega_4 \mathcal{L}_{\mathrm{caus}} .
\end{equation}

As this encoder must be trained first, we pretrain it for a total of $20\,000$ steps
and $20$ epochs (so $1000$ steps per epoch), before freezing its parameters and prepending
it to the actor and the critic. An important parameter is also the dimension of the latent
state, which we set to the actual ``2'' here. Although for later application, one might not know
the real, underlying dimension of the dynamics, we want to make it comparable across the solvers by
using the same low dimension. And as N4SID mostly correctly identifies the true dimension, we use it
here as well. \Cref{alg:encoder-pretraining} summarizes the pretraining.
\begin{algorithm}[htbp]
    \caption{Pretraining of the scalable encoder}
    \label{alg:encoder-pretraining}
    \begin{algorithmic}[1]
        \REQUIRE initial gain $K_0$, action bound $a_{\max}$, encoder $f_\psi \colon \R^{d_{\mathrm{in}}} \to \R^{2}$,
        weights $\omega_1, \dots, \omega_4$
        \FOR{$k = 0, \dots, 20\,000 - 1$}
            \STATE with probability $0.3$, draw $v_k$ uniformly from $[-a_{\max}/10, a_{\max}/10]^{d_{\mathrm{out}}}$,
            otherwise set $v_k = -K_0 y_k$
            \STATE apply $v_k$, store $(y_k, v_k, y_{k+1})$, and reset the plant after $800$ steps
        \ENDFOR
        \FOR{epoch $= 1, \dots, 20$}
            \FOR{each minibatch of $128$ stored transitions}
                \STATE pair every transition with a randomly drawn second transition
                \STATE take an Adam step on $\omega_1 \mathcal{L}_{\mathrm{temp}} + \omega_2 \mathcal{L}_{\mathrm{prop}}
                + \omega_3 \mathcal{L}_{\mathrm{rep}} + \omega_4 \mathcal{L}_{\mathrm{caus}}$ of these pairs
            \ENDFOR
        \ENDFOR
        \STATE standardize the latent state with its mean and standard deviation over the stored observations
        \STATE freeze $\psi$ and place $f_\psi$ in front of the actor and the critic
    \end{algorithmic}
\end{algorithm}

The second trainable extractor we used is based on Todorov et al.'s paper~\cite{Todorov2026Orthogonal}
using orthogonal bottlenecks. It consists of a fully connected MLP, with 512 neurons in the first,
and 256 neurons in the second layer. Following it is a matrix with orthonormal columns, which projects the output
into a state of bottleneck dimension $d_{\mathrm{b}}$, which we here set to $2$. Like $U_1$ in \eqref{eq:lifts}, the matrix
is the orthonormal factor of the QR decomposition of a random matrix with standard normal entries, so that the
feature extractor output $z_k \in \R^{d_{\mathrm{z}}}$ is mapped to
\begin{equation}
    \label{eq:ortho-bottleneck}
    U_{\mathrm{b}}\T z_k \in \R^{d_{\mathrm{b}}} ,
    \qquad U_{\mathrm{b}} \in \R^{d_{\mathrm{z}} \times d_{\mathrm{b}}} , \quad U_{\mathrm{b}}\T U_{\mathrm{b}} = I_{d_{\mathrm{b}}} .
\end{equation}

Different to the previous encoder, this one is trained together with PPO, so we just combine them.

The following figure visualizes the architecture for all our solvers.\needcite{figure below}
\begin{figure}[htbp]
    \centering
    \begin{tikzpicture}[
            x=0.85cm, y=1cm,
            layer/.style={draw, minimum width=0.26cm, inner sep=0pt, line width=0.4pt},
            trained/.style={layer, fill=blue!25},
            frozen/.style={layer, fill=orange!40},
            fixedmap/.style={draw=none, fill=none, font=\scriptsize, inner sep=1pt},
            io/.style={layer, dashed, fill=white},
            link/.style={-{Stealth[length=1.6mm]}, line width=0.4pt},
            size/.style={font=\scriptsize, anchor=north},
            row/.style={font=\small, anchor=east},
        ]
        % columns: input, trunk 1-3, latent, head, actor 1-2, output
        \def\xin{0}\def\xa{1.25}\def\xb{2.35}\def\xc{3.45}\def\xz{4.75}\def\xh{6.0}
        \def\xp{7.25}\def\xq{8.45}\def\xout{9.85}
        % heights: d_in, 512, 256, 64, latent, d_out
        \def\hin{1.5}\def\hfh{1.8}\def\htf{1.45}\def\hsf{0.85}\def\hz{0.3}\def\hout{1.25}

        % Linear PPO
        \def\y{0}
        \node[row] at (-0.5,\y) {Linear PPO};
        \node[io, minimum height=\hin cm] (in) at (\xin,\y) {};
        \node[io, minimum height=\hout cm] (out) at (\xout,\y) {};
        \draw[link] (in) -- node[above, font=\scriptsize] {$W = -K$} (out);

        % Fully-connected PPO
        \def\y{-2.4}\def\lb{0.55}
        \node[row] at (-0.5,\y) {Fully-connected PPO};
        \node[io, minimum height=\hin cm] (in) at (\xin,\y) {};
        \node[trained, minimum height=\hsf cm] (p) at (\xp,\y) {};
        \node[trained, minimum height=\hsf cm] (q) at (\xq,\y) {};
        \node[io, minimum height=\hout cm] (out) at (\xout,\y) {};
        \draw[link] (in) -- (p);
        \draw[link] (p) -- (q);
        \draw[link] (q) -- (out);
        \node[size] at (\xp,\y-\lb) {64};
        \node[size] at (\xq,\y-\lb) {64};

        % State encoder + PPO
        \def\y{-4.8}\def\lb{0.55}
        \node[row] at (-0.5,\y) {State encoder + PPO};
        \node[io, minimum height=\hin cm] (in) at (\xin,\y) {};
        \node[frozen, minimum height=\hz cm] (z) at (\xz,\y) {};
        \node[trained, minimum height=\hsf cm] (p) at (\xp,\y) {};
        \node[trained, minimum height=\hsf cm] (q) at (\xq,\y) {};
        \node[io, minimum height=\hout cm] (out) at (\xout,\y) {};
        \draw[link] (in) -- node[above, font=\scriptsize] {$\hat{C}\T$ (N4SID)} (z);
        \draw[link] (z) -- (p);
        \draw[link] (p) -- (q);
        \draw[link] (q) -- (out);
        \node[size] at (\xz,\y-\lb) {$\hat{n}$};
        \node[size] at (\xp,\y-\lb) {64};
        \node[size] at (\xq,\y-\lb) {64};

        % Scalable encoder + PPO
        \def\y{-7.2}\def\lb{1.0}
        \node[row] at (-0.5,\y) {Scalable encoder + PPO};
        \node[io, minimum height=\hin cm] (in) at (\xin,\y) {};
        \node[frozen, minimum height=\hfh cm] (a) at (\xa,\y) {};
        \node[frozen, minimum height=\htf cm] (b) at (\xb,\y) {};
        \node[frozen, minimum height=\hz cm] (z) at (\xz,\y) {};
        \node[frozen, minimum height=\hsf cm] (h) at (\xh,\y) {};
        \node[trained, minimum height=\hsf cm] (p) at (\xp,\y) {};
        \node[trained, minimum height=\hsf cm] (q) at (\xq,\y) {};
        \node[io, minimum height=\hout cm] (out) at (\xout,\y) {};
        \draw[link] (in) -- (a);
        \draw[link] (a) -- (b);
        \draw[link] (b) -- (z);
        \draw[link] (z) -- (h);
        \draw[link] (h) -- (p);
        \draw[link] (p) -- (q);
        \draw[link] (q) -- (out);
        \node[size] at (\xa,\y-\lb) {512};
        \node[size] at (\xb,\y-\lb) {256};
        \node[size] at (\xz,\y-\lb) {2};
        \node[size] at (\xh,\y-\lb) {64};
        \node[size] at (\xp,\y-\lb) {64};
        \node[size] at (\xq,\y-\lb) {64};

        % Bottleneck PPO
        \def\y{-9.6}\def\lb{1.0}
        \node[row] at (-0.5,\y) {Bottleneck PPO};
        \node[io, minimum height=\hin cm] (in) at (\xin,\y) {};
        \node[trained, minimum height=\hfh cm] (a) at (\xa,\y) {};
        \node[trained, minimum height=\htf cm] (b) at (\xb,\y) {};
        \node[trained, minimum height=\htf cm] (c) at (\xc,\y) {};
        \node[layer, fill=gray!35, minimum height=\hz cm] (z) at (\xz,\y) {};
        \node[trained, minimum height=\hsf cm] (p) at (\xp,\y) {};
        \node[trained, minimum height=\hsf cm] (q) at (\xq,\y) {};
        \node[io, minimum height=\hout cm] (out) at (\xout,\y) {};
        \draw[link] (in) -- (a);
        \draw[link] (a) -- (b);
        \draw[link] (b) -- (c);
        \draw[link] (c) -- node[above, font=\scriptsize] {$U_{\mathrm{b}}\T$} (z);
        \draw[link] (z) -- (p);
        \draw[link] (p) -- (q);
        \draw[link] (q) -- (out);
        \node[size] at (\xa,\y-\lb) {512};
        \node[size] at (\xb,\y-\lb) {256};
        \node[size] at (\xc,\y-\lb) {256};
        \node[size] at (\xz,\y-\lb) {2};
        \node[size] at (\xp,\y-\lb) {64};
        \node[size] at (\xq,\y-\lb) {64};

        % column captions
        \node[font=\scriptsize, align=center] at (\xin,1.3) {observation\\$y_k \in \R^{d_{\mathrm{in}}}$};
        \node[font=\scriptsize, align=center] at (\xb,1.3) {encoder\\(ReLU)};
        \node[font=\scriptsize, align=center] at (\xz,1.3) {latent\\state};
        \node[font=\scriptsize, align=center] at ({(\xp+\xq)/2},1.3) {actor\\(tanh)};
        \node[font=\scriptsize, align=center] at (\xout,1.3) {action\\$v_k \in \R^{d_{\mathrm{out}}}$};

        % legend
        \def\y{-11.3}
        \node[trained, minimum height=0.3cm] (l1) at (-3.6,\y) {};
        \node[right=2pt of l1, font=\scriptsize] (t1) {trained by PPO};
        \node[frozen, minimum height=0.3cm, right=12pt of t1] (l2) {};
        \node[right=2pt of l2, font=\scriptsize] (t2) {fitted before PPO, then frozen};
        \node[layer, fill=gray!35, minimum height=0.3cm, right=12pt of t2] (l3) {};
        \node[right=2pt of l3, font=\scriptsize] {fixed random projection};
    
    \end{tikzpicture}
    \caption{Actor architectures of our PPO variants. For each layer, the number below indicates its neurons. 
    The critic receives the same encoder structure and uses its own 64-64 tanh perceptron with a scalar output.
    Visualization created with Claude Code.}
    \label{fig:ppo-architectures}
\end{figure}

\section{Comparison across Dimensions}
\label{sec:rl-comparison}

Now that we have all our methods set up, we can train and compare them across the different inflation sizes.

\section{Noise on Large Inflation}
\label{sec:rl-noise-large}


\section{Discussion}
\label{sec:rl-discussion}

\cleardoublepage

\chapter{Perception-based Experiments}\label{chap:perception}

\section{Artificial Inflation to Pixel Size}

\subsection{Setting}

\subsection{Results}

\section{The New Environment and Pos-Vel}

\section{Effect Of Visual Inputs On The Solvers}

\subsection{Setting}

\subsection{Results}

\section{Discussion}
\cleardoublepage

\chapter{Conclusion}
\label{chap:conclusion}

% Template

% You generally cover three things in the Conclusions section, and each of these usually merits a separate subsection:

% \begin{enumerate}
% 	\item Conclusions
% 	\item Summary of Contributions
% 	\item Future Research
% \end{enumerate}

% Conclusions are not a rambling summary of the thesis: they are short, concise statements of the inferences that you have made because of your work. It helps to organize these as short numbered paragraphs, ordered from most to least important. All conclusions should be directly related to the research question stated in the introduction. Examples:

% \begin{enumerate}
% 	\item The problem stated in the introduction has been solved: as shown in Sections ? to ??, an algorithm capable of handling large-scale Zylon problems in reasonable time has been developed.
% 	\item The principal mechanism needed in the improved Zylon algorithm is the Grooty mechanism.
% 	\item Etc.
% \end{enumerate}

% The Summary of Contributions will be much sought and carefully read by the examiners. Here you list the contributions of new knowledge that your thesis makes. Of course, the thesis itself must substantiate any claims made here. There is often some overlap with the Conclusions, but that's okay. Concise numbered paragraphs are again best. Organize from most to least important. Examples:

% \begin{enumerate}
% 	\item Developed a much quicker algorithm for large-scale Zylon problems.
% 	\item Demonstrated the first use of the Grooty mechanism for Zylon calculations.
% 	\item Etc.
% \end{enumerate}

% The Future Research subsection is included so that researchers picking up this work in future have the benefit of the ideas that you generated while you were working on the project. Again, concise numbered paragraphs are usually best. 

\section{Conclusions}
\label{sec:conclusions}

In this thesis we have shown that simply enforcing low-dimensional structures or latent layers does not
improve the performance of controllers for overparameterized systems; instead, it might even hurt performance. However,
when combining sufficiently pretrained latent projectors with the right controller architecture, they can
achieve similar performances to fully connected networks, while being more efficient in their parameter count
and training time. While this could be clearly shown for LQR, under RL only tendencies were observed, requiring
further investigation.

We further found conflicting results on whether low-dimensional latent states
can improve robustness to measurement noise. While for LQR solvers and model-informed
\emph{StE} this holds true, for the camera environment, we could not clearly observe that. This however might
be related to the first point, where our RL approaches showed already bad performance
under the clean loop.

Regarding the latent state, it is also important to ensure that not only control information is carried
here, but also information about the observed state. Otherwise, robustness to control noise might worsen.

In relation to that, we also note that a low degradation in control noise can also hint at a barely acting
controller, and is not always a sign of robustness. Still, in our experiments, low-dimensional controllers 
were affected slightly more by it than fully parameterized ones, which we observed across all experiments.

\section{Summary of Contributions}
\label{sec:contributions}

We have presented that combining pretrained encoders with a trainable network
can reduce the parameter count effectively on a toy example, while noting common problems in relation to
architecture design. Failing to translate the success of it onto a vision-based problem brought
additional pitfalls to light. Further, we showed that fully parameterized controllers offer, if even only slightly, more
robustness to control noise, while for the other noise types we could only observe tendencies,
which open possibilities for further research here.

\section{Future Research}
\label{sec:future-research}

From where this thesis ends, there are different directions, which we have left for future research
and discussion. These include, but are not limited to, whether:

\begin{enumerate}
	\item and how the number of parameters can be distributed more effectively between
	the extractor and the actor, to achieve performances comparable to a fully connected network;
	\item our results, such as better robustness to measurement noise of low-dimensional LQR controllers, might hold
	for more complex systems, i.e. a walking humanoid robot;
	\item these findings are transferable to real-world applications, i.e., if this setup
	were built and tested;
\end{enumerate}

Finally, we also want to mention that, while these are interesting directions, the data produced from
our experiments still offers a lot of possibilities for additional analysis. For example,
while we captured the compute time and environment steps needed, a different angle could include
this in the evaluation of the approaches. And while it seems like our conclusions
leave us with more questions than we started with, this is what research is for.
\cleardoublepage

%%%%%%%%%%%%%%%%%%%%%%%%%%%%
% Appendices:
% these are optional! For most Bachelor-theses and some Master-thesis none of them is needed. 
% Just comment them if not needed.
\appendix
% Make plots stick to their headings in the appendix
\let\appendixsection\section
\renewcommand{\section}{\FloatBarrier\appendixsection}
\let\appendixsubsection\subsection
\renewcommand{\subsection}{\FloatBarrier\appendixsubsection}
\fancyhead[LO,RE]{}                      % Define the header style for the appendixpages

\fancyhead[LE,RO]{\it Appendix A. Nomenclature}
  \chapter{Nomenclature}\label{app:nomenclature}

What goes in the appendices? Any material which impedes the smooth development of your presentation, but which is important to justify the results of a thesis. Generally it is material that is of too nitty-gritty a level of detail for inclusion in the main body of the thesis, but which should be available for perusal by the examiners to convince them sufficiently. Examples include program listings, immense tables of data, lengthy mathematical proofs or derivations, etc. 


\section{Notation}
\label{sec:notation}

The following table contains the conventions regarding symbols
used throughout this thesis:

\begin{longtable}{@{}p{0.20\textwidth}p{0.74\textwidth}@{}}
\toprule
\textbf{Symbol} & \textbf{Meaning} \\
\midrule
\endfirsthead
\toprule
\textbf{Symbol} & \textbf{Meaning} \\
\midrule
\endhead
\bottomrule
\endfoot
\multicolumn{2}{@{}l}{\emph{Sets and operators}} \\[2pt]
$I_n$, $\mathbf{1}_n$ & identity matrix and all-ones vector of dimension $n$ \\
$\lambda_i(M)$, $\sigma_i(M)$ & $i$-th eigenvalue and $i$-th singular value of $M$ \\
$\spr(M)$ & spectral radius, $\max_i \lvert \lambda_i(M) \rvert$ \\
$M \succeq 0$, $M \succ 0$ & $M$ symmetric positive semi-definite / positive definite \\
$\lVert \cdot \rVert$, $\lVert \cdot \rVert_F$, $\lVert \cdot \rVert_{*}$ & Euclidean (spectral) norm, Frobenius norm, nuclear norm \\
$\odot$ & elementwise product \\
$\mathcal{N}(\mu, \Sigma)$ & Gaussian distribution with mean $\mu$ and covariance $\Sigma$ \\
$\nabla f$, $\widehat{\nabla} f$ & gradient of $f$ and the estimate of it\\
\midrule
\multicolumn{2}{@{}l}{\emph{System description}} \\[2pt]
$n$, $m$, $q$ & dimension of the state, input and output \\
$t$, $k$, $h$ & continuous time, discrete step index, sampling period \\
$x_k \in \R^{n}$ & state \\
$u_k \in \R^{m}$ & control input computed by the controller \\
$\tilde{u}_k \in \R^{m}$ & control input actually applied to the plant \\
$y_k \in \R^{q}$ & observed output \\
$p$, $v$ & position and velocity of the double integrator, $x = [\,p,\ v\,]\T$ \\
$A_c$, $B_c$ & continuous-time system and input matrix \\
$A$, $B$ & their zero-order-hold discretizations \\
$C$, $D$ & output and feedthrough matrix; $D = 0$ throughout \\
$A_K$ & closed-loop matrix, $A_K = A - BK$ \\
$\mathcal{C}(A,B)$, $\mathcal{O}(C,A)$ & controllability and observability matrix \\
$\hat{x}_k$, $\tilde{x}_k$ & state estimate and estimation error $x_k - \hat{x}_k$ \\
$L$ & observer gain \\
\midrule
\multicolumn{2}{@{}l}{\emph{Optimal control}} \\[2pt]
$Q$, $R$ & state and input cost matrices, $Q \succeq 0$, $R \succ 0$ \\
$K \in \R^{m \times n}$ & feedback gain, $u_k = -K x_k$ \\
$K^{*}$ & optimal gain, solution of \eqref{eq:optimal-gain} \\
$\mathcal{K}$ & set of stabilizing gains, \eqref{eq:stabilizing-set} \\
$J(x_0, u)$ & infinite-horizon cost of an input sequence \\
$P$ & stabilizing solution of the DARE \eqref{eq:dare} \\
$P_K$ & solution of the Lyapunov equation \eqref{eq:lyapunov} for a fixed $K$ \\
$\mathcal{D}$, $\Sigma_0$ & distribution of $x_0$ and its second moment $\E[x_0 x_0\T]$ \\
$f(K)$ & expected cost of a gain, $\tr(P_K \Sigma_0)$ \\
\midrule
\multicolumn{2}{@{}l}{\emph{Noise}} \\[2pt]
$n^{\mathrm{m}}_k$, $\sigma_{\mathrm{m}}$ & measurement noise and its standard deviation \\
$n^{\mathrm{c}}_k$, $\sigma_{\mathrm{c}}$ & control noise (relative) and its standard deviation \\
$\Delta^{A}_k$, $\Delta^{B}_k$, $\sigma_{\mathrm{p}}$ & parametric perturbations and their standard deviation \\
$n^{\mathrm{d}}_k$, $\sigma_{\mathrm{d}}$ & external disturbance and its standard deviation / amplitude \\
\midrule
\multicolumn{2}{@{}l}{\emph{Learning and inflated systems}} \\[2pt]
$d_{in}$, $d_{out}$ & inflated input and output dimension, $d \gg n$ \\
$U_1$, $U_2$ & orthonormal inflation matrices \\
$\lambda$ & weight of the nuclear-norm regulariser \\
$\eta$ & step size \\
$r$ & smoothing radius of the zeroth-order gradient estimator \\
$N$, $T$ & number of rollouts and rollout horizon \\
\end{longtable}

\section{Noise intensities}
\label{sec:app_noise_levels}

The settings used for the different noise intensities, obtained by the
calibration \eqref{eq:noise-calibration} on the non-inflated system.
It is important to note that parametric noise is resampled at every step,
and that the disturbance noise uses its sinusoidal variant \eqref{eq:dist-sin}.

\begin{table}[htbp]
\centering
\caption{Standard deviation per noise channel and intensity, calibrated to the cost ratio $\rho_{\ell}$ of \eqref{eq:noise_calibration}. Every value reaches its ratio within $0.9\,\%$.}
\label{tab:noise_levels}
\begin{tabular}{@{}lrrrrr@{}}
\toprule
Channel & \emph{none} & \emph{weak} & \emph{standard} & \emph{strong} & \emph{divergent} \\
\midrule
Cost ratio $\rho_{\ell}$ & $1$ & $1.2$ & $1.5$ & $3$ & $10$ \\
\midrule
Measurement $\sigma_{\mathrm{m}}$ & $0$ & $0.389$ & $0.607$ & $1.22$ & $2.58$ \\
Control $\sigma_{\mathrm{c}}$ & $0$ & $1.07$ & $1.67$ & $3.19$ & $5.68$ \\
Parametric $\sigma_{\mathrm{p}}$ & $0$ & $0.0211$ & $0.0304$ & $0.0451$ & $0.0582$ \\
Disturbance $\sigma_{\mathrm{d}}$ & $0$ & $0.00767$ & $0.0122$ & $0.0246$ & $0.0519$ \\
\bottomrule
\end{tabular}
\end{table}

\cleardoublepage

\fancyhead[LE,RO]{\it Appendix B. Supplementary Results}
  \chapter{Supplementary Results}
\label{app:supplementary-results}

\section{Reasoning for the performance of TrP under measurement noise}
\label{sec:app-perf-trp-measurment-noise}

% TODO: Apply tildes to all the inflated parts here

We start by looking at the clean loop. Here, any $y_k$ lies in $\operatorname{range}(U_1)$ by construction.
We also have $\tilde{\Sigma}_0 = U_1 \Sigma_0 U_1\T$. % TODO: Why?
Solving now \eqref{eq:state-correlation}, we obtain $\tilde{\Sigma}_{K_j} = U_1 \Sigma_{K_j} U_1\T$, where $\Sigma_{K_j}$ is the solution
of the uninflated loop.

We now define the projector onto the complement of $\operatorname{range}(U_1)$ as $P_{\perp_1} = I_{d_{\mathrm{in}}} - U_1 U_1\T$.
For any $K_j$, we can therefore seperate it into $K_j = K_jU_1 U_1\T + K_jP_{\perp_1}$. If we now take a look at the updated of this second
component in \cref{alg:mb-pg}, we get $K_{j+1}P_{\perp_1} = K_jP_{\perp_1} - \eta \nabla f(K_j)P_{\perp_1}$. Using the gradient calculation
from \eqref{eq:pg-gradient}, the part on the right turns into 
$\nabla f(K_j)P_{\perp_1} = 2 \left( \left( R + B\T P_{K_j} B \right) K_j - B\T P_{K_j} A \right) \tilde{\Sigma}_{K_j} P_{\perp_1}$, where
$\tilde{\Sigma}_{K_j} P_{\perp_1} = U_1 \Sigma_{K_j} U_1\T (I_{d_{\mathrm{in}}} - U_1 U_1\T) = U_1 \Sigma_{K_j} U_1\T - U_1 \Sigma_{K_j} U_1\T = 0$. % TODO: Break into multilines
Plugging this back into \needcite{ref to gradient calc}, we also have $\nabla f(K_j)P_{\perp_1} = 0$, and therefore
$K_{j+1}P_{\perp_1} = K_jP_{\perp_1} - 0$ holds for all $j$, meaning $K_j P_{\perp_1} = K_0 P_{\perp_1}$, i.e. is never updated
via the policy gradient. And as we construct $\tilde{K}_0 = U_2\T K_0 U_1\T + K_{\perp}$, where 
$\tilde{K}_0 P_{\perp_1} = U_2\T K_0 U_1\T P_{\perp_1} + K_{\perp} P_{\perp_1} = 0 + K_{\perp} P_{\perp_1}$, we know this component
to not be zero.

That the gradient equals zero for this part of the gain also makes sense, it is not associated with any commanded actuation, as
$K_{j} P_{\perp_1}(U_1 x) = 0$. However, under measurement noise, where we have the noise component $n_m$, and $y_k = U_1 x + n_{m_k}$

2. \eqref{eq:pg-gradient} multiplies by $\tilde{\Sigma}_K$ from the right, so with $P_{\perp} = I_{d_{\mathrm{in}}} - U_1 U_1\T$ every gradient
satisfies $\nabla f(\tilde{K}_j) P_{\perp} = 0$. Hence $\tilde{K}_j P_{\perp} = \tilde{K}_0 P_{\perp} = K_{\perp} P_{\perp}$ for all $j$
(\eqref{eq:initial-gain}, as $U_1\T P_{\perp} = 0$): \emph{TrP} never changes the part of the gain that sees only the invisible directions.

3. Without noise this part has no effect on the cost. Under measurement noise ($y_k = U_1 x_k + n^{\mathrm{m}}_k$, Ch. 2 l. 544, noise in all
$d_{\mathrm{in}}$ directions), $\tilde{v}_k$ contains $-K_{\perp} P_{\perp} n^{\mathrm{m}}_k$. It is charged through $\tilde{v}_k\T \tilde{R}\,
\tilde{v}_k$ with $\tilde{R} = I$, and enters the dynamics through $B U_2 K_{\perp} P_{\perp} n^{\mathrm{m}}_k$ ($U_2 K_{\perp} U_1 = 0$ does not
imply $U_2 K_{\perp} P_{\perp} = 0$). With $\sigma_{\mathrm{m}}^{2} I_{d_{\mathrm{in}}}$, the control cost alone grows by at least
$\sigma_{\mathrm{m}}^{2} \lVert K_{\perp} P_{\perp} \rVert_F^2$ per step, a term no iteration of \emph{TrP} reduces.

4. S-PI's proximal step soft-thresholds all singular values, so it shrinks this part as well (state as "can remove", unless shown in general).

5. \emph{IdP}: the same argument on the identified model with $\operatorname{range}(\hat{C})$ in place of $\operatorname{range}(U_1)$.


\section{Additional Plots for the Perception Experiments}
\label{sec:app-perc-plots}

\begin{figure}[htbp]
  \centering
  \generatedfigure{ch6_perception}{ch6_perception_trajectories}
  \caption[Example trajectories of the PPO variants.]{Trajectories of the PPO variants for one sampled episode under one seed.}
  \label{fig:app-perc-trajectories}
\end{figure}


\section{Exit Rates of the Perception Experiments under Noise}
\label{sec:app-perc-noise-terminated}

\Cref{tab:perc-noise-terminated} notes how often the mass leaves the screen under each noise type and
intensity in \cref{subsec:perception-noise-results}.

\begin{table}[htbp]
\centering
\caption{Share of test episodes in the camera environment that end with the mass leaving the view, in percent, per noise channel and intensity. Mean $\pm$ standard deviation over the seeds; \emph{none} is the clean loop.}
\label{tab:perc-noise-terminated}
\begin{tabular}{@{}llrrrrr@{}}
\toprule
Channel & Solver & \emph{none} & \emph{weak} & \emph{standard} & \emph{strong} & \emph{divergent} \\
\midrule
Measurement & FcE & $0$ & $0 \pm 1$ & $3 \pm 2$ & $14 \pm 11$ & $66 \pm 26$ \\
 & OrB & $0$ & $0$ & $3 \pm 7$ & $19 \pm 27$ & $66 \pm 25$ \\
 & ScE & $6 \pm 3$ & $8 \pm 3$ & $10 \pm 5$ & $17 \pm 9$ & $42 \pm 30$ \\
 & PVE & $3 \pm 8$ & $18 \pm 24$ & $43 \pm 41$ & $73 \pm 34$ & $79 \pm 27$ \\
\midrule
Control & FcE & $0$ & $0$ & $0$ & $0$ & $0$ \\
 & OrB & $0$ & $0$ & $0$ & $0$ & $0$ \\
 & ScE & $6 \pm 3$ & $6 \pm 3$ & $7 \pm 4$ & $10 \pm 5$ & $18 \pm 5$ \\
 & PVE & $3 \pm 8$ & $4 \pm 9$ & $4 \pm 9$ & $11 \pm 18$ & $21 \pm 28$ \\
\midrule
Parametric & FcE & $0$ & $8 \pm 3$ & $8 \pm 3$ & $16 \pm 2$ & $42 \pm 5$ \\
 & OrB & $0$ & $8 \pm 3$ & $9 \pm 2$ & $16 \pm 5$ & $39 \pm 5$ \\
 & ScE & $6 \pm 3$ & $52 \pm 6$ & $52 \pm 6$ & $62 \pm 5$ & $74 \pm 3$ \\
 & PVE & $3 \pm 8$ & $36 \pm 22$ & $37 \pm 23$ & $48 \pm 20$ & $62 \pm 13$ \\
\midrule
Disturbance & FcE & $0$ & $0$ & $0$ & $0 \pm 1$ & $89 \pm 3$ \\
 & OrB & $0$ & $0$ & $0$ & $5 \pm 3$ & $85 \pm 5$ \\
 & ScE & $6 \pm 3$ & $9 \pm 4$ & $10 \pm 5$ & $18 \pm 5$ & $46 \pm 7$ \\
 & PVE & $3 \pm 8$ & $3 \pm 8$ & $4 \pm 8$ & $8 \pm 10$ & $36 \pm 15$ \\
\bottomrule
\end{tabular}
\end{table}


\cleardoublepage

\fancyhead[LE,RO]{\it Appendix C. Calibration Results}
  \chapter{Complete Simulation Results}\label{app:completeResults}

\begin{table}[htbp]
\centering
\caption{S-PI (model-free) with $\lambda = 400$ under \emph{medium} inflation, over the seeds. Mean $\pm$ standard deviation over 10 runs.}
\label{tab:appendix-spi-mf-medium-400}
\resizebox{\ifdim\width>\linewidth\linewidth\else\width\fi}{!}{%
\begin{tabular}{lrrrrr}
\toprule
Step size & 3.5e-09 & 3.5e-08 & 3.5e-07 & 3.5e-06 & 3.5e-05 \\
\midrule
True cost $J$ & $1.204 \times 10^{4} \pm 416.6$ & $9080 \pm 42.58$ & $7527 \pm 37.58$ & $6374 \pm 31.73$ & $1.241 \times 10^{4} \pm 1.917 \times 10^{-12}$ \\
Training suboptimality $(J-J^*)/J^*$ & $1.06 \pm 0.0712$ & $0.553 \pm 0.00728$ & $0.287 \pm 0.00643$ & $0.09 \pm 0.00543$ & $1.12 \pm 2.96 \times 10^{-16}$ \\
Gain rank & $25.0 \pm 0.0$ & $25.0 \pm 0.0$ & $25.0 \pm 0.0$ & $25.0 \pm 0.0$ & $25.0 \pm 0.0$ \\
Spectral radius & $0.9975 \pm 0.0000$ & $0.9975 \pm 0.0000$ & $0.9974 \pm 0.0000$ & $0.9959 \pm 0.0001$ & $0.9975 \pm 0.0000$ \\
Iterations & $50 \pm 0$ & $50 \pm 0$ & $50 \pm 0$ & $50 \pm 0$ & $50 \pm 0$ \\
Env.\ steps to convergence & $8 \times 10^{6} \pm 0$ & $8 \times 10^{6} \pm 0$ & $8 \times 10^{6} \pm 0$ & $8 \times 10^{6} \pm 0$ & $6.47 \times 10^{6} \pm 3.23 \times 10^{6}$ \\
Env.\ steps to best policy & $3.47 \times 10^{6} \pm 3.71 \times 10^{6}$ & $7.98 \times 10^{6} \pm 5.06 \times 10^{4}$ & $7.98 \times 10^{6} \pm 5.06 \times 10^{4}$ & $7.98 \times 10^{6} \pm 5.06 \times 10^{4}$ & $1.6 \times 10^{5} \pm 0$ \\
\bottomrule
\end{tabular}
}
\end{table}


\begin{table}[htbp]
\centering
\caption{S-PI (model-free) with $\lambda = 4000$ under \emph{medium} inflation, over the seeds. Mean $\pm$ standard deviation over 10 runs.}
\label{tab:appendix-spi-mf-medium-4000}
\resizebox{\ifdim\width>\linewidth\linewidth\else\width\fi}{!}{%
\begin{tabular}{lrrrrr}
\toprule
Step size & 3.5e-09 & 3.5e-08 & 3.5e-07 & 3.5e-06 & 3.5e-05 \\
\midrule
True cost $J$ & $1.196 \times 10^{4} \pm 417.9$ & $9108 \pm 45.74$ & $7839 \pm 45.92$ & $7238 \pm 85.8$ & $1.241 \times 10^{4} \pm 2.908 \times 10^{-12}$ \\
Training suboptimality $(J-J^*)/J^*$ & $1.04 \pm 0.0715$ & $0.557 \pm 0.00782$ & $0.34 \pm 0.00785$ & $0.238 \pm 0.0147$ & $1.12 \pm 4.68 \times 10^{-16}$ \\
Gain rank & $25.0 \pm 0.0$ & $25.0 \pm 0.0$ & $22.5 \pm 1.4$ & $9.7 \pm 2.9$ & $25.0 \pm 0.0$ \\
Spectral radius & $0.9975 \pm 0.0000$ & $0.9976 \pm 0.0000$ & $0.9977 \pm 0.0000$ & $0.9973 \pm 0.0001$ & $0.9975 \pm 0.0000$ \\
Iterations & $50 \pm 0$ & $50 \pm 0$ & $50 \pm 0$ & $50 \pm 0$ & $50 \pm 0$ \\
Env.\ steps to convergence & $8 \times 10^{6} \pm 0$ & $8 \times 10^{6} \pm 0$ & $8 \times 10^{6} \pm 0$ & $8 \times 10^{6} \pm 0$ & $6.47 \times 10^{6} \pm 3.23 \times 10^{6}$ \\
Env.\ steps to best policy & $4.26 \times 10^{6} \pm 3.76 \times 10^{6}$ & $8 \times 10^{6} \pm 0$ & $8 \times 10^{6} \pm 0$ & $8 \times 10^{6} \pm 0$ & $1.6 \times 10^{5} \pm 0$ \\
\bottomrule
\end{tabular}
}
\end{table}


\begin{table}[htbp]
\centering
\caption{S-PI (model-free) with $\lambda = 40000$ under \emph{medium} inflation, over the seeds. Mean $\pm$ standard deviation over 10 runs.}
\label{tab:appendix-spi-mf-medium-40000}
\resizebox{\ifdim\width>\linewidth\linewidth\else\width\fi}{!}{%
\begin{tabular}{lrrrrr}
\toprule
Step size & 3.5e-09 & 3.5e-08 & 3.5e-07 & 3.5e-06 & 3.5e-05 \\
\midrule
True cost $J$ & $1.16 \times 10^{4} \pm 29.37$ & $9466 \pm 40.8$ & $2.996 \times 10^{4} \pm 2887$ & $2.978 \times 10^{4} \pm 8845$ & $1.241 \times 10^{4} \pm 1.917 \times 10^{-12}$ \\
Training suboptimality $(J-J^*)/J^*$ & $0.983 \pm 0.00502$ & $0.619 \pm 0.00698$ & $4.12 \pm 0.494$ & $4.09 \pm 1.51$ & $1.12 \pm 2.96 \times 10^{-16}$ \\
Gain rank & $25.0 \pm 0.0$ & $17.9 \pm 0.7$ & $15.5 \pm 2.4$ & $7.3 \pm 5.0$ & $25.0 \pm 0.0$ \\
Spectral radius & $0.9975 \pm 0.0000$ & $0.9979 \pm 0.0000$ & $0.9996 \pm 0.0000$ & $0.9996 \pm 0.0001$ & $0.9975 \pm 0.0000$ \\
Iterations & $50 \pm 0$ & $50 \pm 0$ & $50 \pm 0$ & $50 \pm 0$ & $50 \pm 0$ \\
Env.\ steps to convergence & $8 \times 10^{6} \pm 0$ & $8 \times 10^{6} \pm 0$ & $8 \times 10^{6} \pm 0$ & $8 \times 10^{6} \pm 0$ & $6.18 \times 10^{6} \pm 3.08 \times 10^{6}$ \\
Env.\ steps to best policy & $8 \times 10^{6} \pm 0$ & $8 \times 10^{6} \pm 0$ & $8 \times 10^{6} \pm 0$ & $1.14 \times 10^{6} \pm 2.19 \times 10^{5}$ & $1.6 \times 10^{5} \pm 0$ \\
\bottomrule
\end{tabular}
}
\end{table}


\begin{table}[htbp]
\centering
\caption{S-PI (model-free) with $\lambda = 400000$ under \emph{medium} inflation, over the seeds. Mean $\pm$ standard deviation over 10 runs. Stable counts the seeds whose controller left the plant stable.}
\label{tab:appendix-spi-mf-medium-400000}
\resizebox{\ifdim\width>\linewidth\linewidth\else\width\fi}{!}{%
\begin{tabular}{lrrrrr}
\toprule
Step size & 3.5e-09 & 3.5e-08 & 3.5e-07 & 3.5e-06 & 3.5e-05 \\
\midrule
Stable & $10/10$ & $9/10$ & $4/10$ & $0/10$ & $0/10$ \\
\midrule
True cost $J$ & $1.205 \times 10^{4} \pm 60.32$ & $1.139 \times 10^{6} \pm 5.329 \times 10^{5}$ & $2.785 \times 10^{9} \pm 5.205 \times 10^{9}$ & \multicolumn{1}{c}{--} & \multicolumn{1}{c}{--} \\
Training suboptimality $(J-J^*)/J^*$ & $1.06 \pm 0.0103$ & $194 \pm 91.1$ & $4.76 \times 10^{5} \pm 8.9 \times 10^{5}$ & \multicolumn{1}{c}{--} & \multicolumn{1}{c}{--} \\
Gain rank & $17.9 \pm 0.6$ & $8.0 \pm 3.1$ & $1.0 \pm 0.0$ & \multicolumn{1}{c}{--} & \multicolumn{1}{c}{--} \\
Spectral radius & $0.9978 \pm 0.0000$ & $1.0000 \pm 0.0000$ & $1.0000 \pm 0.0000$ & \multicolumn{1}{c}{--} & \multicolumn{1}{c}{--} \\
Iterations & $50 \pm 0$ & $50 \pm 0$ & $50 \pm 0$ & \multicolumn{1}{c}{--} & \multicolumn{1}{c}{--} \\
Env.\ steps to convergence & $8 \times 10^{6} \pm 0$ & $8 \times 10^{6} \pm 0$ & $8 \times 10^{6} \pm 0$ & \multicolumn{1}{c}{--} & \multicolumn{1}{c}{--} \\
Env.\ steps to best policy & $8 \times 10^{6} \pm 0$ & $8 \times 10^{6} \pm 0$ & $2.2 \times 10^{6} \pm 1.85 \times 10^{6}$ & \multicolumn{1}{c}{--} & \multicolumn{1}{c}{--} \\
\bottomrule
\end{tabular}
}
\end{table}


\begin{table}[htbp]
\centering
\resizebox{\ifdim\width>\linewidth\linewidth\else\width\fi}{!}{%
\begin{tabular}{lrrrrr}
\toprule
Step size & 3.5e-09 & 3.5e-08 & 3.5e-07 & 3.5e-06 & 3.5e-05 \\
\midrule
Stable & $0/10$ & $0/10$ & $0/10$ & $0/10$ & $0/10$ \\
\midrule
True cost $J$ & \multicolumn{1}{c}{--} & \multicolumn{1}{c}{--} & \multicolumn{1}{c}{--} & \multicolumn{1}{c}{--} & \multicolumn{1}{c}{--} \\
Training suboptimality $(J-J^*)/J^*$ & \multicolumn{1}{c}{--} & \multicolumn{1}{c}{--} & \multicolumn{1}{c}{--} & \multicolumn{1}{c}{--} & \multicolumn{1}{c}{--} \\
Gain rank & \multicolumn{1}{c}{--} & \multicolumn{1}{c}{--} & \multicolumn{1}{c}{--} & \multicolumn{1}{c}{--} & \multicolumn{1}{c}{--} \\
Spectral radius & \multicolumn{1}{c}{--} & \multicolumn{1}{c}{--} & \multicolumn{1}{c}{--} & \multicolumn{1}{c}{--} & \multicolumn{1}{c}{--} \\
Test cost & \multicolumn{1}{c}{--} & \multicolumn{1}{c}{--} & \multicolumn{1}{c}{--} & \multicolumn{1}{c}{--} & \multicolumn{1}{c}{--} \\
\bottomrule
\end{tabular}
}
\caption{Model-free S-PI at lambda=4e+06, medium inflation, over 10 seeds. Stable counts the seeds whose controller left the plant stable.}
\label{tab:appendix-spimf-medium-4e06}
\end{table}

\cleardoublepage

% ... add as much appendices as you need (one can also add source code, for example)

%\fancyhead[LE]{\it \leftmark}
%\chapter{}
\fancyhead[LE,RO]{\it Bibliography}       % A bibliography never have a letter or numbering!
    \bibliographystyle{plain}             % Style for presenting the literature
    \addcontentsline{toc}{chapter}{Bibliography}% Add to the TOC
    \bibliography{thesis}
\cleardoublepage

%%%%%%%%%%%%%%%%%%%%%%%%%%%%

% Formal page 1
\vspace{2cm}
\chapter*{Verwendete Hilfsmittel}
\label{sec:tool-usage}
\fancyhead[LE]{\it Erkl\"arung zur Benutzung von Hilfsmitteln}
Für die Erstellung dieser Bachelorarbeit wurde Generative KI innerhalb
der im Folgenden genannten Hilfsmittel für die angeführten Zwecke verwendet. Der Autor
hat jegliche produzierten Ausgaben ausführlich geprüft, und übernimmt Verantwortung für diese.
Die Selbständigkeit der produzierten Ideen, sowie deren Umsetzung, waren zu jeder
Zeit gewährleistet.

Im Detail wurde genutzt:

GitHub Copilot, integriert in Visual Studio Code, für:
\begin{itemize}
  \item Automatische Vervollständigung von Codezeilen ("Inline Suggestions") zur Unterstützung
  bei Latex Formatierung, sowie der python-basierten Experimente
\end{itemize}

Claude Code, mit den Modellen Opus 4.8, Opus 5 und Opus 5.5, für:
\begin{itemize}
  \item Korrektur von Grammatik, Zeichensetzung und Tippfehlern
  \item Formatierung von Latex Code
  \item Quellenvorschläge
  \item Automatisches Einpflegen meiner Quellen in die Latex Bibliografie
  \item Erstellung von Latex-basierten Schema Visualisierungen von Code Implementierungen
  \item Unterstützung bei der Erstellung von Skripten für die automatische Generierung von Figuren und Tabellen basierend
  auf den Experiment Ergebnissen
  \item Erstellung einer Anleitung, sowie eines Skripts für vereinfachte Datentransfers zwischen dem TUHH HPC Cluster und
  meinem lokalen Gerät
  \item Anmerkung von Inkonsistenzen und Logik Fehlern in der textlichen Ausarbeitung
  \item Herausfordern meiner Standpunkte und Umsetzung ("Sparring")
  \item Erkennen von Fehlern im Python Code
\end{itemize}

%The backcover is always empty
\newpage
\thispagestyle{empty}
\hspace{1cm}
\newpage

%%%%%%%%%%%%%%%%%%%%%%%%%%%%
% Formal page 2
\vspace{2cm}
\chapter*{Erkl\"arung der Urheberschaft}
\label{sec:urheber}
\fancyhead[LE]{\it Erkl\"arung der Urheberschaft}
Ich versichere an Eides statt, dass ich die \trtype{} im Studiengang \trstudies{} selbstst\"andig verfasst und keine anderen als die angegebenen Hilfsmittel -- insbesondere keine im Quellenverzeichnis nicht benannten Internet-Quellen -- benutzt habe. Es wurden keine Quellen genutzt, die nicht genannt wurden. Alle Stellen, die w\"ortlich oder sinngem\"a{\ss} aus Ver\"offentlichungen entnommen wurden, sind als solche kenntlich gemacht. Ich versichere weiterhin, dass ich die Arbeit vorher nicht in einem anderen Pr\"ufungsverfahren eingereicht habe und die eingereichte schriftliche Fassung der auf dem elektronischen Speichermedium entspricht.

%Ich versichere an Eides statt, dass ich die vorliegende \trtype{} selbstst\"andig und ohne unerlaubte Hilfe Dritter angefertigt habe. Alle Stellen, die inhaltlich oder w\"ortlich aus anderen Ver\"offentlichungen stammen, sind kenntlich gemacht. Diese Arbeit lag in gleicher oder \"ahnlicher Weise noch keiner Pr\"ufungsbeh\"orde vor und wurde bisher noch nicht ver\"offentlicht.

\vspace{4cm}
\noindent Ort, Datum \hfill Unterschrift

%The backcover is always empty
\newpage
\thispagestyle{empty}
\hspace{1cm}
\newpage

%%%%%%%%%%%%%%%%%%%%%%%%%%%%
% Formal page 3
\vspace{2cm}
\chapter*{Erkl\"arung zur Ver\"offentlichung}
\label{sec:veroeffentlichung}
\fancyhead[LE]{\it Erkl\"arung zur Ver\"offentlichung}
Ich erkl\"are mein Einverst\"andnis mit der Einstellung dieser \trtype{} in den Bestand der Bibliothek.

\vspace{4cm}
\noindent Ort, Datum \hfill Unterschrift

%The backcover is always empty
\newpage
\thispagestyle{empty}
\hspace{1cm}
\newpage

\end{document}
%EOF

